%% 第24回編集委員会（物理第3巻）：組 波W4
%% p3-08.tex  波-8　重なる波：干渉・回折・ドップラー効果（★2〜4・3段にしない）
%% 中身は計画書 tools/phys3-plan.md の B「波-8」。公式の持ち主は A（A18・A27〜A32・A44・A45）。
%%   うなり（A18）は 4-9⑦の計算を指し、物理の場面で導き直す。単スリット（A29）は 9-4④の方形パルスの変換を指す。
%%   マッハ数は熱-4 ⑦（熱-4 は音速を a と書く。この巻は c）。マッハ円すいはこの章⑥（熱-4 ⑨が予告している）。
%%   光への反射・屈折の当てはめは波-12 ③、光の干渉・回折は波-13（ヤングの実験・回折格子・円い穴）。
%%   ④の単スリットの最初の暗い所 a sinθ＝λ は、波-13 の組の頼み（まとめ役から）で、式と導き方をはっきり書いた。
%% 数値はすべて Python で計算した：checks/物理第3巻/波-8/calc_p3_08.py（出力 calc_p3_08_out.txt）、
%%   hamen_p3_08.py（⑦のプログラム。本文の phcode と同じ）、kanji_p3_07_08.py（常用漢字の確かめ）
%% 音速：本文の値は 20 ℃の 343 m/s（熱-4 ⑥）。救急車のサイレン（はてな？・例題1・7・11）だけは、計画書のとおり「340 m/s とする」と断り、
%%   343 m/s とのちがい（1 Hz 以下）を⑤に書いた。
%% 記号（この章）：音の速さ（波の速さ）c、音源の速さ v_s・観測者の速さ v_o（音源→観測者の向きを正）、振動数 f・聞こえる振動数 f′、
%%   波長 λ、干渉の次数 m、スリットの幅 a、波源の間かく d、道のり l₁・l₂、マッハ数 Ma（熱-4 と同じ立体の2文字）、円すいの半角 α、
%%   風の速さ w、斜めの角 θ。高校の教科書の V（音速）・v（波の速さ）は「思い出そう」で一言。
%% 出典（第24回の追記：書くときに確かめた。11月5日。WebSearch で一次資料を引く資料を2つ以上照らし合わせた）
%%   ・救急車のサイレン 960 Hz・770 Hz：東京消防庁の資料 https://www.tfd.metro.tokyo.lg.jp/content/safetyreport/000007279.pdf（計画書 D）
%%   ・サイレンの大きさ（前方20 m で 90 dB 以上 120 dB 以下）：道路運送車両の保安基準の細目を定める告示 第231条（国土交通省 https://www.mlit.go.jp/jidosha/content/S231.pdf を curl と pdftotext で読んだ）
%%   ・20 ℃の空気の音速 343 m/s：熱-4 ⑥。CoolProp 8.0（空気の状態式）で 343.3 m/s・1.2046 kg/m³ を確かめた
%%   ・速度測定のレーダー：総務省の無線局の審査基準に「10GHz帯又は24GHz帯…速度測定」の区分がある
%%       （https://www.soumu.go.jp/main_content/000240215.pdf）。24.15 GHz・10.525 GHz は米国の K・X バンドの値
%%       （copradar.com、hypertextbook.com）。本文は「24 GHz 帯（ここでは 24.15 GHz とする）」と書いた
%%   ・体の中の音速 1540 m/s：radiopaedia "Acoustic impedance"、Stanford RAD225 の講義資料
%%   ・アブラコウモリの探索の声 約40 kHz：日本心理学会の資料（psych.or.jp/wp-content/uploads/2020/04/89-27-28.pdf）など2つ
%%   ・キクガシラコウモリのドップラーシフト補償：H.-U. Schnitzler, Z. vergl. Physiol. 57, 376–408 (1968)、同 J. comp. Physiol. 82, 79–92 (1973)（書誌は検索で見つけた論文の引用で確かめた。本文は読んでいない）
%%   ・ドップラー 1842年（"Über das farbige Licht der Doppelsterne…"、5月25日発表）、ボイス・バロット 1845年（ユトレヒト近くの列車のラッパ）：
%%       Buys Ballot, Annalen der Physik und Chemie 66, 321–351 (1845)（Zenodo record 1423606）、Doppler, Abh. Königl. Böhm. Ges. Wiss. (5) 2, 465–482 (1842)、Taylor & Francis の章
%%   ・ハッブル定数：Planck 2018 で 67.4±0.5、SH0ES（Riess ら 2022）で 73.04±1.04 km/(s·Mpc)、約5σ：arXiv:2203.06142（レビュー）、CERN Courier
%%   ・船の引き波の角 19.47°（arcsin 1/3）：Rabaud & Moisy, PRL 110, 214503 (2013)、arXiv:1304.2653 "Ship wakes: Kelvin or Mach angle?"（速いボートでは小さくなる）
%%   ・ベル X-1（1947年10月14日、イェーガー、マッハ約1.06・時速約1126 km）：Britannica "X-1 airplane"、スミソニアン航空宇宙博物館、space.com（Mach 1.05〜1.06 と資料で少し違う）
%%   ・ホイヘンスの原理（1678年にパリの科学アカデミーで発表、1690年『光についての論考』）
%% 〔要確認〕（確かめられなかった値）
%%   ・踏切の警報音 700 Hz は「とする」で使った（実際の値は調べていない）
%%   ・校種：計画書 A の案のまま（うなり＝高校物理基礎、ホイヘンス・反射屈折・干渉・回折・ドップラー＝高校物理、
%%     斜めのドップラー・衝撃波＝高校物理（発展））。
%%     整合の組が、高等学校学習指導要領（平成30年告示）解説 理科編（文部科学省
%%     https://www.mext.go.jp/content/20250311-mxt_kyoiku02-100002620_05.pdf）で照らした：「物理」(2)(ア)(イ) に
%%     ホイヘンスの原理・反射・屈折・正弦波の式・干渉と回折・音のドップラー効果（観測者と音源が同一直線上を動く場合を中心に）。
%%     「音源が音速以上の速さで動いているときに起こる現象に触れることも考えられる」。「物理基礎」に共振・共鳴・うなり。
%%     斜めのドップラー・衝撃波の角を（発展）とするのと合う。
\makeatletter
% 救急車（横から）  \rxicar{(後ろの下)}{大きさ}
\newcommand\rxicar[2]{\begin{scope}[shift={#1}, scale=#2]
  \filldraw[fill=white, draw=bkInk!80, line width=0.6pt, rounded corners=0.8mm] (0,1.2) rectangle (14,7.5);
  \filldraw[fill=white, draw=bkInk!80, line width=0.6pt, rounded corners=0.8mm] (14,1.2) -- (19,1.2) -- (19,4.8) -- (16.5,7.5) -- (14,7.5) -- cycle;
  \fill[bkBlue!25] (15,4.6) -- (16.8,6.8) -- (15,6.8) -- cycle;
  \fill[bkRed!80] (0.3,3.6) rectangle (13.7,4.4);
  \fill[bkRed] (12,7.5) rectangle (14,8.6);
  \fill[bkInk!85] (3.5,1.2) circle (1.4) (15.5,1.2) circle (1.4);\end{scope}}
% 動く音源の波面（円をずらして描く）  \rxiwaves{比 v/c}{間かく}{数}{色}
\newcommand\rxiwaves[4]{\foreach \n in {1,...,#3} {\draw[#4, line width=0.6pt] ({-#1*#2*\n},0) circle ({#2*\n});}}
\makeatother

\phchapter{8}{重なる波：\\干渉・回折・ドップラー効果}{救急車のサイレンは、なぜ通りすぎると低く聞こえる？}{\bklevelrange{2}{4}}{重なる波：干渉・回折・ドップラー効果}
\bkmichishirube{4-9,5-17,9-4,熱-4,波-5,波-7}{波-9,波-10,波-12,波-13}

\begin{goal}
\begin{itemize}[label={\color{bkTheme}●}, leftmargin=1.6\zw, itemsep=0pt, topsep=1pt]
\item \bklevel{2}\enspace 近い高さの2つの音が重なるときの、うなりの回数を求められる。
\item \bklevel{3}\enspace ホイヘンスの原理で、反射・屈折の法則と回折を説明できる。
\item \bklevel{3}\enspace 2つの波源からの道のりの差で、強め合う所と弱め合う所を求められる。
\item \bklevel{3}\enspace 波面を数えてドップラー効果の式を導き、聞こえる高さを求められる。
\item \bklevel{4}\enspace 斜めのドップラー効果・反射板・衝撃波で、速さをはかるしくみを説明できる。
\end{itemize}
\end{goal}

\phkoushiki{%
\phk{うなりの回数 $|f_1-f_2|$}{高校物理基礎}{ph:波-8:unari}
\phk{ホイヘンスの原理、反射の法則、屈折の法則 $\sin\theta_1/\sin\theta_2=c_1/c_2=\lambda_1/\lambda_2$}{高校物理}{ph:波-8:kussetsu}
\phk{2つの波源の干渉 $|l_1-l_2|=m\lambda$、$(m+\tfrac12)\lambda$}{高校物理}{ph:波-8:kanshou}
\phk{回折（すきまが波長くらいで大きく回りこむ）}{高校物理}{ph:波-8:kaisetsu}
\phk{ドップラー効果 $f'=f\,(c-v_{\mathrm o})/(c-v_{\mathrm s})$（風・反射板があるときも）}{高校物理}{ph:波-8:doppler}
\phk{斜めのドップラー効果（$v\cos\theta$）}{高校物理（発展）}{ph:波-8:naname}
\phk{衝撃波の角 $\sin\alpha=c/v_{\mathrm s}$}{高校物理（発展）}{ph:波-8:mach}
}

\begin{hatena}
\noindent\begin{minipage}[c]{0.52\linewidth}\raggedright
「ピーポー、ピーポー」と、救急車が近づいてきます。目の前を通りすぎたとたん、サイレンの音が、ストンと低くなります。救急車は、通りすぎたからといって、サイレンを切りかえているわけではありません。運転している人には、ずっと同じ高さの音が聞こえています。

外で聞いている人にだけ、音の高さが変わって聞こえるのは、なぜでしょう。
\end{minipage}\hfill
\begin{minipage}[c]{0.45\linewidth}\centering
\begin{tikzpicture}[x=0.9mm, y=0.9mm]
  \begin{scope}
  \clip (-30,-16) rectangle (30,18);
  \rxiwaves{0.4}{5}{8}{bkBlue!70}
  \end{scope}
  \rxicar{(-6,-2.5)}{0.55}
  \draw[bkRed!80!black, line width=1pt, -{Stealth[length=2.2mm]}] (6,-6) -- (14,-6);
  \node[font=\sffamily\footnotesize, anchor=south, fill=white, inner sep=1pt] at (20,13) {前：つまる};
  \node[font=\sffamily\footnotesize, anchor=south, fill=white, inner sep=1pt] at (-19,13) {後ろ：のびる};
  \node[font=\sffamily\footnotesize, anchor=north, fill=white, inner sep=1pt] at (0,-16.5) {進む向き →};
\end{tikzpicture}
\end{minipage}
\end{hatena}

\noindent 波-7までは、1つの波が進み、はね返るようすを見てきました。この章では、2つ以上の波が重なると何が起きるか（\textbf{干渉}・\textbf{うなり}）、波がすきまや物のかげに回りこむこと（\textbf{回折}）、そして、音を出す物や聞く人が動くと高さが変わること（\textbf{ドップラー効果}）を調べます。どれも、波の山（波面）を図にかいて、数えれば説明できます。

\begin{yattemiyou}[2つの音を重ねる・動く音を聞く]
\textbf{用意するもの}：スマホ2台（決まった高さの音を出すアプリを入れる）\par\vspace{2pt}
\begin{enumerate}[label={\bkcircled[bkGreen]{\arabic*}}, leftmargin=2.4\zw, itemsep=1pt, topsep=0pt]
\item 1台で\qty{440}{Hz}、もう1台で\qty{440}{Hz}の音を、同時に鳴らす。次に、片方を\qty{442}{Hz}、\qty{444}{Hz}に変える。音の大きさが「ワーン、ワーン」とゆれる回数を、10秒間数える。
\item 2台を\qty{1}{m}ほどはなして同じ\qty{1000}{Hz}の音を鳴らし、片耳をふさいで、2台の前をゆっくり横に歩く。音が大きい所と小さい所があるか調べる。
\item 車の通りの多い道の、安全な歩道に立つ。前を通りすぎる車やバイクの音が、通りすぎる前と後でどう変わるか聞く（道路には出ない）。
\end{enumerate}
\end{yattemiyou}

\noindent やってみようの答え：①\qty{440}{Hz}どうしでは、ゆれは起きません。\qty{442}{Hz}にすると1秒に約2回、\qty{444}{Hz}では約4回、音の大きさがゆれます。ゆれの回数は、2つの高さの差になっています（②で説明する\textbf{うなり}）。②歩くと、音が大きく聞こえる所と、小さく聞こえる所が、何十 cm かおきに並んでいます（③の\textbf{干渉}）。③車の音は、近づくあいだは高く、通りすぎると急に低くなります（⑤の\textbf{ドップラー効果}）。

\begin{omoidasou}
\bkrecall{波-5}{波の速さ $c$、振動数 $f$、波長 $\lambda$ の間に $c=f\lambda$。波は重ねても、そのまま通りぬける（重ね合わせ：波-5 ⑥）。}
\bkrecall{4-9}{$\sin A+\sin B=2\sin\dfrac{A+B}{2}\cos\dfrac{A-B}{2}$。\qty{440}{Hz}と\qty{442}{Hz}の音を重ねると、1秒に2回うなる。}
\bkrecall{9-4}{はば2の方形パルスのフーリエ変換は $2\sin\omega/\omega$ で、はじめて0になるのは $\omega=\pi$。はばを半分にすると、広がりは2倍。}
\bkrecall{5-17}{$x$ が小さいとき、$1/(1-x)\approx1+x$（1次近似）。}
\bkrecall{熱-4}{\qty{20}{\celsius}の空気の音の速さは\qty{343}{m/s}。マッハ数 $\mathrm{Ma}$＝（速さ）÷（音の速さ）。熱-4 では音の速さを $a$ と書いた。}
\end{omoidasou}

\noindent この章の記号：高校の教科書は、音の速さを $V$、波の速さを $v$ と書きます。この巻では、どちらも $c$ と書きます（波-5）。音を出す物（\textbf{音源}）の速さを $v_{\mathrm s}$、聞く人（\textbf{観測者}）の速さを $v_{\mathrm o}$ と書きます（s は source、o は observer の頭文字）。

%% ======================== 物理の道具 ========================
\bksection{doguu}{物理の道具}

\bksubsection{1}{ホイヘンスの原理と、反射・屈折の法則}
\label{ph:波-8:sec1}
池に石を落とすと、輪の形の波が広がります。波の山（または谷）をつないだ線（面）を\textbf{波面}といいます。\bkidx{はめん}{波面}小さな音源からの波面は球面、遠くから来る波や、長い板をゆすって作る波の波面は、平らな面になります（平面波）。波の進む向きは、波面にいつも直角です。

\phboxhead{ホイヘンスの原理}
次の時刻の波面は、どうやって決まるのでしょう。オランダのホイヘンス（1678年ごろ）は、次のように考えました。
\begin{bkcenter}
波面の上のどの点も、新しい小さな波（\textbf{素元波}）の源になる。\\
素元波は、波の速さ $c$ で球面に広がる。$\Delta t$ 後の波面は、\\
たくさんの素元波に共通に接する面（包絡面）になる。
\end{bkcenter}
\label{ph:波-8:huygens}これを\textbf{ホイヘンスの原理}といいます。\bkidx{ほいへんすのげんり}{ホイヘンスの原理}

\begin{bkfigure}
\begin{tikzpicture}[x=1mm, y=1mm]
  % 平面波
  \draw[bkBlue!80!black, line width=1.2pt] (0,0) -- (0,36);
  \foreach \yy in {3,9,15,21,27,33} {\fill[bkRed] (0,\yy) circle (0.8); \draw[bkGray, line width=0.5pt] (0,\yy) ++(-90:8) arc (-90:90:8);}
  \draw[bkBlue!80!black, line width=1.2pt, dashed] (8,0) -- (8,36);
  \draw[bkOrange!85!black, line width=1pt, -{Stealth[length=2.2mm]}] (12,18) -- (22,18);
  \node[font=\sffamily\footnotesize, anchor=north] at (4,-9) {平面波};
  \draw[bkInk, line width=0.5pt, {Stealth[length=2pt]}-{Stealth[length=2pt]}] (0,41) -- (8,41);
  \node[font=\sffamily\footnotesize, anchor=south] at (4,42) {$c\,\Delta t$};
  % 球面波
  \begin{scope}[shift={(60,18)}]
  \fill[bkYellow!70!black] (0,0) circle (1.2);
  \draw[bkBlue!80!black, line width=1.2pt] (0,0) circle (10);
  \foreach \a in {0,30,...,330} {\fill[bkRed] (\a:10) circle (0.8); \draw[bkGray, line width=0.5pt] (\a:10) circle (7);}
  \draw[bkBlue!80!black, line width=1.2pt, dashed] (0,0) circle (17);
  \node[font=\sffamily\footnotesize, anchor=north] at (0,-19.5) {球面波};
  \end{scope}
  \node[anchor=west, font=\small, text width=36mm, align=left] at (84,18) {%
  実線：今の波面\par
  ●：素元波の源\par
  細い円：$\Delta t$ のあいだに広がった素元波\par
  破線：$\Delta t$ 後の波面（包絡面）};
\end{tikzpicture}
\bkfigcaption{ホイヘンスの原理。平面波は平面のまま、球面波は球面のまま進む。}\label{ph:波-8:fig-huygens}
\end{bkfigure}

\phboxhead{反射の法則}
平面波が、平らな壁に斜めに当たります。波面と壁の交わる点Aに波面が着いた時刻から、波面の反対のはしBが壁のB$'$に着くまでの時間を $\Delta t$ とすると、$\mathrm{BB'}=c\,\Delta t$ です。そのあいだに、Aから出た素元波も、半径 $c\,\Delta t$ に広がっています。反射した波面は、B$'$からこの素元波に引いた接線A$'$B$'$です。

\begin{bkfigure}
\begin{tikzpicture}[x=1mm, y=1mm]
  \fill[bkGray!25] (-4,-5) rectangle (64,0);
  \draw[bkInk, line width=0.8pt] (-4,0) -- (64,0);
  \coordinate (A) at (10,0); \coordinate (Bp) at (50,0);
  \coordinate (B) at ({10+40*cos(40)*cos(40)},{40*cos(40)*sin(40)});
  \coordinate (Ap) at ({50-40*cos(40)*cos(40)},{40*cos(40)*sin(40)});
  \draw[bkBlue!80!black, line width=1.2pt] (A) -- (B);
  \draw[bkRed!80!black, line width=1.2pt] (Ap) -- (Bp);
  \draw[bkGray, line width=0.5pt] (A) ++(0:{40*sin(40)}) arc (0:180:{40*sin(40)});
  \draw[bkInk, line width=0.5pt, dashed] (B) -- (Bp);
  \draw[bkInk, line width=0.5pt, dashed] (A) -- (Ap);
  \fill (A) circle (0.8) (B) circle (0.8) (Ap) circle (0.8) (Bp) circle (0.8);
  \node[font=\sffamily\footnotesize, anchor=north] at (A) {A};
  \node[font=\sffamily\footnotesize, anchor=north] at ([yshift=-0.3mm]Bp) {B$'$};
  \node[font=\sffamily\footnotesize, anchor=south west] at (B) {B};
  \node[font=\sffamily\footnotesize, anchor=south east] at (Ap) {A$'$};
  \node[font=\sffamily\footnotesize, anchor=west, fill=white, inner sep=1pt] at ($(B)!0.55!(Bp)+(1.5,0)$) {$c\,\Delta t$};
  \node[font=\sffamily\footnotesize, anchor=east, fill=white, inner sep=1pt] at ($(A)!0.55!(Ap)+(-1.5,0)$) {$c\,\Delta t$};
  \node[anchor=west, font=\small, text width=50mm, align=left] at (70,16) {%
  直角三角形 ABB$'$ と B$'$A$'$A は、\par 斜辺 AB$'$ が共通で、\par BB$'$＝AA$'$＝$c\,\Delta t$。\par\vspace{2pt}
  だから合同で、来る角と\par はね返る角は等しい。};
\end{tikzpicture}
\bkfigcaption{ホイヘンスの原理で作る反射の法則。来た波面 AB が壁に当たり、$\Delta t$ 後には、はね返った波面 A$'$B$'$ になる。}\label{ph:波-8:fig-hansha}
\end{bkfigure}

\noindent 2つの直角三角形 ABB$'$ と B$'$A$'$A は、斜辺 AB$'$ が共通で、もう1辺が等しい（どちらも $c\,\Delta t$）ので、合同です。波面が壁となす角が、来る波とはね返る波で等しくなります。波の進む向きは波面に直角なので、壁に立てた垂線と進む向きのなす角についても、
\begin{bkcenter}
\textbf{入射角＝反射角}（反射の法則）
\end{bkcenter}
が成り立ちます。光の反射（波-12 ②）も、山びこの音の反射も、同じ法則に従います。

\phboxhead{屈折の法則}
こんどは、波の速さが $c_1$ の所（媒質1）から、$c_2$ の所（媒質2）へ、波が斜めに入ります。波面の点Aが境目に着いてから、点BがB$'$に着くまでの時間を $\Delta t$ とすると、$\mathrm{BB'}=c_1\Delta t$。そのあいだに、Aから媒質2の中へ出た素元波の半径は $c_2\Delta t$ です。

\begin{bkfigure}
\begin{tikzpicture}[x=1mm, y=1mm]
  \fill[bkBlue!8] (-4,-34) rectangle (64,0);
  \draw[bkInk, line width=0.8pt] (-4,0) -- (64,0);
  \node[font=\sffamily\footnotesize, anchor=north west] at (-3,34) {媒質1（速さ $c_1$）};
  \node[font=\sffamily\footnotesize, anchor=south west] at (-3,-33) {媒質2（速さ $c_2$、遅い）};
  \coordinate (A) at (10,0); \coordinate (Bp) at (50,0);
  \coordinate (B) at ({10+40*cos(50)*cos(50)},{40*cos(50)*sin(50)});
  \draw[bkBlue!80!black, line width=1.2pt] (A) -- (B);
  \draw[bkInk, line width=0.5pt, dashed] (B) -- (Bp);
  % 素元波の半径 c2Δt = 40 sin50 × (c2/c1)，c2/c1=0.6
  \pgfmathsetmacro\rr{40*sin(50)*0.6}
  \draw[bkGray, line width=0.5pt] (A) ++(-180:\rr) arc (-180:0:\rr);
  \pgfmathsetmacro\ang{asin(\rr/40)}
  \coordinate (Ap) at ($(A)+({-90+\ang}:{\rr})$);
  \coordinate (Ap) at ($(A)!{sin(\ang)*sin(\ang)}!(Bp)+(0,{-\rr*cos(\ang)})$);
  \draw[bkRed!80!black, line width=1.2pt] (Ap) -- (Bp);
  \draw[bkInk, line width=0.5pt, dashed] (A) -- (Ap);
  \fill (A) circle (0.8) (B) circle (0.8) (Ap) circle (0.8) (Bp) circle (0.8);
  \node[font=\sffamily\footnotesize, anchor=south] at ([yshift=0.4mm]A) {A};
  \node[font=\sffamily\footnotesize, anchor=south] at ([yshift=0.4mm]Bp) {B$'$};
  \node[font=\sffamily\footnotesize, anchor=south east] at (B) {B};
  \node[font=\sffamily\footnotesize, anchor=north] at (Ap) {A$'$};
  \node[font=\sffamily\footnotesize, anchor=west, fill=white, inner sep=1pt] at ($(B)!0.5!(Bp)+(1,0)$) {$c_1\Delta t$};
  \node[font=\sffamily\footnotesize, anchor=east, fill=white, inner sep=1pt] at ($(A)!0.5!(Ap)+(-1,0)$) {$c_2\Delta t$};
  \node[anchor=west, font=\small, text width=50mm, align=left] at (70,0) {%
  $\sin\theta_1=\dfrac{\mathrm{BB'}}{\mathrm{AB'}}=\dfrac{c_1\Delta t}{\mathrm{AB'}}$\par\vspace{3pt}
  $\sin\theta_2=\dfrac{\mathrm{AA'}}{\mathrm{AB'}}=\dfrac{c_2\Delta t}{\mathrm{AB'}}$\par\vspace{3pt}
  $\theta_1$：入射角（波面と境目の角＝\par 進む向きと垂線の角）\par
  $\theta_2$：屈折角};
\end{tikzpicture}
\bkfigcaption{ホイヘンスの原理で作る屈折の法則（$c_2<c_1$ のとき）。遅い媒質2では、素元波が小さいので、波面の向きが変わる。}\label{ph:波-8:fig-kussetsu}
\end{bkfigure}

\noindent 屈折した波面は、B$'$から素元波に引いた接線A$'$B$'$です。直角三角形 ABB$'$ と AA$'$B$'$ で、斜辺 AB$'$ は共通です。入射角 $\theta_1$ と屈折角 $\theta_2$ の $\sin$ は、図\ref{ph:波-8:fig-kussetsu}の右の式なので、2つをわって AB$'$ と $\Delta t$ を消すと、
\pheq{\frac{\sin\theta_1}{\sin\theta_2}=\frac{c_1}{c_2}=\frac{\lambda_1}{\lambda_2}=n_{12}}{単位なし}
\label{ph:波-8:kussetsu}となります。これが\textbf{屈折の法則}です。\bkidx{くっせつのほうそく}{屈折の法則}$n_{12}$ を、媒質1に対する媒質2の\textbf{屈折率}といいます。2つめの等号は、境目の両側で振動数 $f$ が同じだから成り立ちます（境目の点は、来た波にゆすられるリズムで、媒質2をゆするから。$\lambda=c/f$ なので、$\lambda$ は $c$ に比例）。

\begin{phflow}[ホイヘンスの原理から]
\phflowrow{ホイヘンスの原理（素元波の包絡面）}{同じ媒質の中で\\はね返る}{入射角＝反射角}
\phflowrow{ホイヘンスの原理（素元波の包絡面）}{速さの違う媒質へ\\振動数は変わらない}{$\dfrac{\sin\theta_1}{\sin\theta_2}=\dfrac{c_1}{c_2}=\dfrac{\lambda_1}{\lambda_2}$}
\end{phflow}

\begin{phyomi}[$\dfrac{\sin\theta_1}{\sin\theta_2}=\dfrac{c_1}{c_2}$]
「遅い媒質へ入ると、波は境目の垂線に近づくように曲がる」。$c_2<c_1$ なら $\sin\theta_2<\sin\theta_1$ で、屈折角は入射角より小さくなります。速い媒質へ入ると、垂線から離れるように曲がります。光でも同じで、光が空気から水へ入ると垂線に近づきます（波-12 ③）。
\end{phyomi}

\begin{phmijika}[夜は、遠くの電車の音がよく聞こえる]
晴れた夜は、地面が冷え、地面に近い空気ほど冷たく、上の空気のほうが暖かくなることがあります。音は暖かい空気の中ほど速いので（熱-4 ⑥）、上へ向かった音は、速い層に斜めに入って垂線から離れる向きに曲がり、少しずつ下へ曲がって、地面にもどってきます。昼間は逆に、地面の近くが暖かいので、音は上へ曲がって逃げていきます。夜に、昼は聞こえない遠くの電車や工場の音が聞こえるのはこのためです。
\end{phmijika}

\begin{seikaku}
ホイヘンスの原理をそのまま使うと、素元波は後ろ向きにも広がるので、後ろへ進む波もできてしまいそうです。実際には後ろへの波はできません。のちにフレネルとキルヒホッフが、素元波の強さが向きによって変わることを、波動方程式から示して、この問題を解きました（専門の本で）。この本では、包絡面の作図として使います。
\end{seikaku}

\bksubsection{2}{うなり：近い高さの音を重ねる}
\label{ph:波-8:sec2}
振幅が同じ $A$ で、振動数が少し違う $f_1$ と $f_2$ の音（$f_1>f_2$）を、同じ場所で重ねます（重ね合わせ：波-5 ⑥）。耳の所での空気のゆれは、
\[y=A\sin2\pi f_1t+A\sin2\pi f_2t\]
です。4-9⑦の和を積に直す式 $\sin A+\sin B=2\sin\tfrac{A+B}{2}\cos\tfrac{A-B}{2}$ を使うと、
\[y=\underbrace{2A\cos\Bigl(2\pi\,\frac{f_1-f_2}{2}\,t\Bigr)}_{\text{ゆっくり変わる振幅}}\times\underbrace{\sin\Bigl(2\pi\,\frac{f_1+f_2}{2}\,t\Bigr)}_{\text{まん中の高さの音}}\]
となります。わたしたちには、2つのまん中の高さ $(f_1+f_2)/2$ の音が、ゆっくり大きくなったり小さくなったりして聞こえます。

音の大きさは、振幅 $2A\cos(\cdots)$ の\textbf{大きさ}で決まります。$\cos$ は1周期のあいだに $+1$ と $-1$ の2回、大きさが最大になるので、音が大きくなる回数は、1秒に $\cos$ の振動数 $(f_1-f_2)/2$ の2倍です。
\pheq{\text{1秒あたりのうなりの回数}=|f_1-f_2|}{Hz ＝ Hz－Hz}
\label{ph:波-8:unari}これを\textbf{うなり}といいます。\bkidx{うなり}{うなり}

\begin{bkfigure}
\begin{tikzpicture}
\begin{axis}[bkplot, axis lines=left, width=124mm, height=40mm, xmin=0, xmax=1.08, ymin=-2.5, ymax=2.7, xtick={0.25,0.5,0.75,1}, ytick={-2,2}, xlabel={$t$（s）}, xticklabel style={/pgf/number format/fixed}]
  \addplot[bkBlue, line width=0.8pt, domain=0:1, samples=700] {sin(360*22*x)+sin(360*20*x)};
  \addplot[bkOrange, line width=1pt, dash pattern=on 3pt off 1.5pt, domain=0:1, samples=120] {2*cos(360*x)};
  \addplot[bkOrange, line width=1pt, dash pattern=on 3pt off 1.5pt, domain=0:1, samples=120] {-2*cos(360*x)};
\end{axis}
\end{tikzpicture}
\bkfigcaption{\qty{22}{Hz}と\qty{20}{Hz}の波の和（実線。見やすいように低い振動数でかいた）と、振幅 $\pm2\cos(2\pi\times1\times t)$（破線）。音が大きくなるのは、0秒・0.5秒・1秒の、1秒に2回（$22-20=2$）。}\label{ph:波-8:fig-unari}
\end{bkfigure}

\begin{phbaisuru}{振動数の差 $f_1-f_2$}
うなりの回数も2倍になります。差が小さいほど、うなりはゆっくりになり、差が0になると、うなりは消えます。差が1秒に10回をこえるくらいになると、ゆれとしては聞こえず、「ざらざらした」音に聞こえ、もっと差が大きいと、2つの別々の高さの音に聞こえます。
\end{phbaisuru}

\begin{phmijika}[楽器の調律]
ギターやバイオリンの弦を合わせるとき、基準の音（おんさやチューナーの\qty{440}{Hz}）と弦の音を同時に鳴らし、うなりがゆっくりになるように弦の張力を変えます。うなりが消えたら、高さがぴったり合ったしるしです。プロペラが2つある飛行機で「ウォンウォン」という音がするのも、2つのエンジンの回転が少し違うときのうなりです。
\end{phmijika}

\noindent うなりの「ふくらみ」が進む速さが、波-5 ⑦で学んだ群速度です（波-7 ⑥）。

\bksubsection{3}{2つの波源の干渉}
\label{ph:波-8:sec3}
同じ振動数・同じ振幅で、山と谷を同時に出す（\textbf{同じ位相の}）2つの波源 S$_1$・S$_2$ を考えます。2つのスピーカーに同じ音を入れたときや、水面を同時にたたく2本の棒です。ある点Pまでの道のりを $l_1$・$l_2$ とします。

点Pでは、S$_1$ からの波 $A\sin(\omega t-kl_1)$ と、S$_2$ からの波 $A\sin(\omega t-kl_2)$ が重なります（ここでは、道のりで振幅が変わることを無視する）。和を積に直す式（4-9⑦）で、
\[y=2A\cos\Bigl(\frac{k(l_1-l_2)}{2}\Bigr)\sin\Bigl(\omega t-\frac{k(l_1+l_2)}{2}\Bigr)\]
です。点Pのゆれの振幅は $2A\bigl|\cos\bigl(\pi\,(l_1-l_2)/\lambda\bigr)\bigr|$（$k=2\pi/\lambda$）で、\textbf{道のりの差} $l_1-l_2$ だけで決まります。
\pheq{|l_1-l_2|=m\lambda\ \text{で強め合い、}\ |l_1-l_2|=\Bigl(m+\frac12\Bigr)\lambda\ \text{で弱め合う}\quad(m=0,1,2,\cdots)}{m ＝ m}
\label{ph:波-8:kanshou}道のりの差が波長のちょうど整数倍なら、山と山・谷と谷が重なって振幅は $2A$、半波長ずれていれば、山と谷が重なって打ち消し合い、振幅は0です。このように、波が重なって強め合ったり弱め合ったりすることを\textbf{干渉}といいます。\bkidx{かんしょう}{干渉}$m$ を干渉の\textbf{次数}といいます。

\begin{bkfigure}
\begin{tikzpicture}[x=1mm, y=1mm]
  \begin{scope}
  \clip (-38,0) rectangle (38,44);
  \foreach \r in {6,12,...,54} {\draw[bkBlue!60, line width=0.6pt] (-9,0) circle (\r); \draw[bkBlue!60, line width=0.6pt] (9,0) circle (\r);}
  \foreach \r in {3,9,...,51} {\draw[bkOrange!70, line width=0.5pt, dotted] (-9,0) circle (\r); \draw[bkOrange!70, line width=0.5pt, dotted] (9,0) circle (\r);}
  % 強め合う線（双曲線）：差 0, ±6, ±12
  \draw[bkRed!80!black, line width=1.4pt] (0,0) -- (0,44);
  \foreach \m in {1,2} {\pgfmathsetmacro\aa{3*\m}\pgfmathsetmacro\bb{sqrt(81-\aa*\aa)}
    \draw[bkRed!80!black, line width=1.4pt, domain=0:2.2, samples=40, smooth] plot ({\aa*cosh(\x)},{\bb*sinh(\x)});
    \draw[bkRed!80!black, line width=1.4pt, domain=0:2.2, samples=40, smooth] plot ({-\aa*cosh(\x)},{\bb*sinh(\x)});}
  % 弱め合う線：差 ±3, ±9
  \foreach \m in {0,1} {\pgfmathsetmacro\aa{1.5+3*\m}\pgfmathsetmacro\bb{sqrt(81-\aa*\aa)}
    \draw[bkInk, line width=1.0pt, dashed, domain=0:2.2, samples=40, smooth] plot ({\aa*cosh(\x)},{\bb*sinh(\x)});
    \draw[bkInk, line width=1.0pt, dashed, domain=0:2.2, samples=40, smooth] plot ({-\aa*cosh(\x)},{\bb*sinh(\x)});}
  \end{scope}
  \fill[bkInk] (-9,0) circle (1.2) (9,0) circle (1.2);
  \node[font=\sffamily\footnotesize, anchor=north] at (-9,-1.5) {S$_1$};
  \node[font=\sffamily\footnotesize, anchor=north] at (9,-1.5) {S$_2$};
  \node[anchor=west, font=\small, text width=50mm, align=left] at (41,24) {%
  細い実線：山、点線：谷（波長 $\lambda$ は線6\,mm分）\par\vspace{2pt}
  太い実線：強め合う線（差 $0$、$\pm\lambda$、$\pm2\lambda$）\par\vspace{2pt}
  破線：弱め合う線（差 $\pm\tfrac12\lambda$、$\pm\tfrac32\lambda$）\par\vspace{2pt}
  どの線も、2点からの距離の差が一定の曲線（双曲線）。};
\end{tikzpicture}
\bkfigcaption{2つの波源の干渉（波源の間かくは $3\lambda$）。山と山が重なる点を結ぶと強め合う線、山と谷が重なる点を結ぶと弱め合う線になる。}\label{ph:波-8:fig-kanshou}
\end{bkfigure}

\begin{phbaisuru}[2倍]{波長 $\lambda$（振動数を半分）}
強め合う線と弱め合う線の間かくが広がり、線の数は半分ほどになります。低い音ほど、大きい所と小さい所の間が広くなります。
\end{phbaisuru}

\begin{phmijika}[音を消すヘッドホン]
雑音を消すヘッドホン（ノイズキャンセリング）は、小さなマイクで周りの雑音を聞き、それと山と谷が逆になった音（位相が $\pi$ ずれた音）を、耳元のスピーカーから出します。耳の所で、雑音と逆の音が干渉して打ち消し合うので、静かになります。ステレオの2つのスピーカーの線を、片方だけ＋と－を逆につなぐと、低い音が打ち消し合って、音がやせて聞こえるのも同じ理由です。
\end{phmijika}

\bksubsection{4}{回折：波は回りこむ}
\label{ph:波-8:sec4}
塀の向こうにいる人の声は聞こえますが、姿は見えません。波は、すきまを通りぬけたり、物のはしを通ったりするとき、物のかげへ回りこみます。これを\textbf{回折}といいます。\bkidx{かいせつ}{回折}\label{ph:波-8:kaisetsu}ホイヘンスの原理で考えると、すきまの中の各点が素元波の源になり、素元波はかげの向きへも広がるからです。回りこみ方は、すきまの幅 $a$ と波長 $\lambda$ の比で決まります（図\ref{ph:波-8:fig-kaisetsu}）。
\begin{itemize}[leftmargin=1.2\zw, itemsep=1pt, topsep=2pt, label={\color{bkTheme}●}]
\item $\lambda$ が $a$ と同じくらいか、それより長い：すきまを出た波は、ほとんどすべての向きへ広がる。
\item $\lambda$ が $a$ よりずっと短い：波はほとんどまっすぐ進み、少しだけ広がる。
\end{itemize}

\begin{bkfigure}
\begin{tikzpicture}[x=1mm, y=1mm]
  % 左：広いすきま（a は波長の4倍）
  \begin{scope}
  \fill[bkGray!50] (0,-24) rectangle (2,-8) (0,8) rectangle (2,24);
  \foreach \xx in {-16,-12,-8,-4} {\draw[bkBlue!80!black, line width=0.8pt] (\xx,-22) -- (\xx,22);}
  \foreach \n in {1,...,6} {\pgfmathsetmacro\rr{4*\n}
    \draw[bkBlue!80!black, line width=0.8pt] (2+\rr,-8) -- (2+\rr,8);
    \draw[bkBlue!80!black, line width=0.8pt] (2+\rr,8) arc (0:50:\rr);
    \draw[bkBlue!80!black, line width=0.8pt] (2+\rr,-8) arc (0:-50:\rr);}
  \node[font=\sffamily\footnotesize, anchor=north] at (10,-25) {広いすきま（$a=4\lambda$）};
  \end{scope}
  % 右：せまいすきま（a は波長の半分）
  \begin{scope}[shift={(66,0)}]
  \fill[bkGray!50] (0,-24) rectangle (2,-1) (0,1) rectangle (2,24);
  \foreach \xx in {-16,-12,-8,-4} {\draw[bkBlue!80!black, line width=0.8pt] (\xx,-22) -- (\xx,22);}
  \foreach \n in {1,...,6} {\pgfmathsetmacro\rr{4*\n} \draw[bkBlue!80!black, line width=0.8pt] (2,0) ++(-85:\rr) arc (-85:85:\rr);}
  \node[font=\sffamily\footnotesize, anchor=north] at (10,-25) {せまいすきま（$a=\tfrac12\lambda$）};
  \end{scope}
\end{tikzpicture}
\bkfigcaption{すきまを通る波（線は波の山。間かくが波長）。広いすきまでは、波はほぼまっすぐ進み、はしの近くだけ少し回りこむ。せまいすきまでは、すきまが1つの素元波の源になり、半円の形に広がる。}\label{ph:波-8:fig-kaisetsu}
\end{bkfigure}

\phboxhead{単スリット：最初の暗い所 $a\sin\theta=\lambda$}
幅 $a$ のすきま（\textbf{スリット}）に、平面波がまっすぐ当たります。すきまの中の点が出す素元波が、遠くで、もとの進む向きから角 $\theta$ の向きへ進むとき、すきまの上のはしと下のはしからの道のりの差は $a\sin\theta$ です。

\begin{bkfigure}
\begin{tikzpicture}[x=1.5mm, y=1.5mm]
  \fill[bkGray!50] (0,13) rectangle (2,24) (0,-5) rectangle (2,7);
  \foreach \xx in {-10,-7.5,-5} {\draw[bkBlue!80!black, line width=1pt] (\xx,8) -- (\xx,12);}
  \draw[bkOrange!85!black, line width=1pt, -{Stealth[length=2mm]}] (-11,16) -- (-4,16);
  \node[font=\sffamily\footnotesize, anchor=south] at (-7.5,16.5) {来る波};
  \foreach \yy in {13,10,7} {\fill[bkRed] (2,\yy) circle (0.45); \draw[bkInk, line width=0.6pt, -{Stealth[length=1.8mm]}] (2,\yy) -- ++(30:20);}
  \draw[bkInk, line width=0.5pt, dashed] (2,13) -- (4.6,8.5);
  \draw[bkRed!80!black, line width=1.6pt] (2,7) -- (4.6,8.5);
  \node[font=\sffamily\footnotesize, anchor=north west, text=bkRed!80!black, inner sep=1pt] at (3.6,6.6) {$a\sin\theta$};
  \draw[bkInk, line width=0.5pt] (2,7) -- (20,7);
  \draw[bkInk, line width=0.5pt] (14,7) arc (0:30:12);
  \node[font=\sffamily\footnotesize] at (16,9.2) {$\theta$};
  \draw[bkInk, line width=0.5pt, {Stealth[length=2pt]}-{Stealth[length=2pt]}] (-2,7) -- (-2,13);
  \node[font=\sffamily\footnotesize, anchor=east] at (-2.3,10) {$a$};
  \node[anchor=west, font=\small, text width=60mm, align=left] at (23,9) {%
  すきまを上半分と下半分に分ける。上半分の1つの点と、そこから $\tfrac a2$ 下の点（下半分の対応する点）の道のりの差は $\tfrac a2\sin\theta$。\par\vspace{2pt}
  これが $\tfrac{\lambda}{2}$ なら、2つは打ち消し合う。上半分のすべての点に相手があるので、全体で0。\par\vspace{2pt}
  だから、最初の暗い所は $a\sin\theta=\lambda$。};
\end{tikzpicture}
\bkfigcaption{単スリットの回折。上のはしと下のはしからの道のりの差は $a\sin\theta$（赤の短い線）。}\label{ph:波-8:fig-slit}
\end{bkfigure}

\noindent すきまを上半分と下半分に分けて、上半分の点と、ちょうど $\tfrac a2$ 下の点を組にします。組の2点の道のりの差は $\tfrac a2\sin\theta$ です。これが半波長 $\tfrac\lambda2$ なら、組の2つの素元波は打ち消し合います（③）。上半分のどの点にも相手があるので、すきま全体からの波が打ち消し合い、その向きには波が来ません。だから、
\pheq{a\sin\theta=\lambda\qquad\text{（最初の暗い所）}}{m ＝ m}
\label{ph:波-8:slit}です。まっすぐの向き（$\theta=0$）では、すべての素元波が同じ道のりで強め合い、いちばん明るく（大きく）なります。その両側の、$\sin\theta=\lambda/a$ の向きまでの範囲に、波の大部分が広がって進みます。

この結果は、9-4④の方形パルスの変換と同じ形です。すきまの中の「1が並んでいる区間」を、時間のパルスの代わりに場所のパルスと見ると、遠くでの波の振幅は $\sin(\pi a\sin\theta/\lambda)\big/(\pi a\sin\theta/\lambda)$ に比例し、9-4④の $2\sin\omega/\omega$ と同じ形の曲線になります。はじめて0になるのは $\pi a\sin\theta/\lambda=\pi$、つまり $a\sin\theta=\lambda$ で、上の答えと一致します。9-4で「時間で短いほど、周波数で広い」と学んだことが、ここでは「すきまがせまいほど、広く回りこむ」になっています。

\begin{bkfigure}
\begin{tikzpicture}
\begin{axis}[bkplot, axis lines=left, width=110mm, height=40mm, xmin=-3.3, xmax=3.3, ymin=0, ymax=1.18, xtick={-3,-2,-1,1,2,3}, xticklabels={$-3$,$-2$,$-1$,1,2,3}, ytick={1}, xlabel={$\dfrac{a\sin\theta}{\lambda}$}, ylabel={強さ}]
  \addplot[bkcurve, domain=-3.2:3.2, samples=400] {x==0 ? 1 : (sin(deg(pi*x))/(pi*x))^2};
  \node[font=\sffamily\footnotesize, anchor=south west, fill=white, inner sep=1pt] at (axis cs:1.05,0.12) {最初の暗い所};
\end{axis}
\end{tikzpicture}
\bkfigcaption{単スリットの回折の強さ（まっすぐの向きを1とした）。$a\sin\theta=\pm\lambda$ で0になり、その外側に弱い山がならぶ。}\label{ph:波-8:fig-slitI}
\end{bkfigure}

\begin{phbaisuru}[半分]{すきまの幅 $a$}
$\sin\theta=\lambda/a$ が2倍になるので、角が小さいときは、広がる角がほぼ2倍になります。すきまをせまくするほど、波は広がります。同じすきまなら、波長の長い波ほどよく広がります。
\end{phbaisuru}

\begin{phmijika}[声は聞こえて、姿は見えない]
人の声の波長は、\qty{500}{Hz}で約\qty{0.7}{m}。戸口や塀のはしの大きさ（\qty{1}{m}くらい）と同じくらいなので、声は大きく回りこんで、かげにいる人にも届きます。光の波長は約 $0.5\times10^{-6}$\,m と、戸口よりはるかに短いので、ほとんど回りこまず、くっきりした影ができます（例題5）。隣の部屋の音楽で、低い音（波長が長い）ばかりがよく聞こえるのも、回折の働きです。
\end{phmijika}

\bksubsection{5}{ドップラー効果：波面を数える}
\label{ph:波-8:sec5}
音を出す物や、聞く人が動くと、聞こえる音の高さ（振動数）が変わります。これを\textbf{ドップラー効果}といいます。\bkidx{どっぷらーこうか}{ドップラー効果}音源と観測者は1本の直線の上にあり、空気は止まっている（風がない）とします。速さは、\textbf{音源から観測者へ向かう向きを正}とします。

\phboxhead{音源が止まっているとき}
音源は1秒に $f$ 個の波面（山）を出します。波面は、出た所から、空気に対して速さ $c$ で広がります。1秒後には、はじめに出た波面は $c$〔m〕先にあり、そのあいだの長さ $c$ に $f$ 個の波面がならんでいます。波面の間かく（波長）は $\lambda=c/f$ です。

\phboxhead{音源が動くとき：波面の間かくが変わる}
音源が、観測者に向かって速さ $v_{\mathrm s}$ で動いているとします。1秒のあいだに、はじめに出た波面は $c$ 進みますが、音源も $v_{\mathrm s}$ 進みながら次々に波面を出すので、前の $f$ 個の波面は、長さ $c-v_{\mathrm s}$ の中につめこまれます（図\ref{ph:波-8:fig-doppler}）。前の波長は
\[\lambda'=\frac{c-v_{\mathrm s}}{f}\]
と短くなります。後ろでは、長さ $c+v_{\mathrm s}$ に広がって、波長は $(c+v_{\mathrm s})/f$ と長くなります。

\begin{bkfigure}
\begin{tikzpicture}[x=0.95mm, y=0.95mm]
  \begin{scope}
  \clip (-62,-22) rectangle (40,22);
  \rxiwaves{0.35}{6}{12}{bkBlue!70}
  \end{scope}
  \rxicar{(-5.5,-2.3)}{0.55}
  \draw[bkRed!80!black, line width=1pt, -{Stealth[length=2.2mm]}] (6,-5) -- (14,-5);
  \node[font=\sffamily\footnotesize, anchor=north, fill=white, inner sep=1pt] at (10,-5.5) {$v_{\mathrm s}$};
  \fill[bkGreen!60!black] (36,0) circle (1.5);
  \node[font=\sffamily\footnotesize, anchor=south, fill=white, inner sep=1pt] at (36,2.2) {前の人};
  \fill[bkGreen!60!black] (-58,0) circle (1.5);
  \node[font=\sffamily\footnotesize, anchor=south, fill=white, inner sep=1pt] at (-56,2.2) {後ろの人};
  \node[font=\sffamily\footnotesize, fill=white, inner sep=1pt, anchor=north] at (24,-17) {せまい：$\dfrac{c-v_{\mathrm s}}{f}$};
  \node[font=\sffamily\footnotesize, fill=white, inner sep=1pt, anchor=north] at (-40,-17) {広い：$\dfrac{c+v_{\mathrm s}}{f}$};
\end{tikzpicture}
\bkfigcaption{動く音源の波面（音源の速さを音の速さの0.35倍にして、大げさにかいた）。どの円も、それを出したときの音源の位置が中心。前では波面がつまり、後ろではのびる。}\label{ph:波-8:fig-doppler}
\end{bkfigure}

\phboxhead{観測者が聞く振動数}
止まっている観測者には、波面が速さ $c$ でやって来ます。1秒に耳を通る波面の数（聞こえる振動数 $f'$）は、$c$ を波面の間かく $\lambda'$ でわった数です。観測者も、音源から遠ざかる向きに速さ $v_{\mathrm o}$ で動いているときは、波面が観測者を追いこす速さが $c-v_{\mathrm o}$ になるので、
\[f'=\frac{c-v_{\mathrm o}}{\lambda'}\]
です。$\lambda'=(c-v_{\mathrm s})/f$ を入れると、
\pheq{f'=\frac{c-v_{\mathrm o}}{c-v_{\mathrm s}}\,f}{Hz ＝ $\dfrac{\text{m/s}}{\text{m/s}}$×Hz}
\label{ph:波-8:doppler}$v_{\mathrm s}$・$v_{\mathrm o}$ は、音源から観測者へ向かう向きを正とした速さです。音源が観測者に近づくと $v_{\mathrm s}>0$ で分母が小さくなって $f'$ は高く、観測者が音源に近づくと $v_{\mathrm o}<0$ で分子が大きくなって、やはり高くなります。

\phboxhead{はてな？の答え}
\label{ph:波-8:hatena}\bkhatenaref の答えです。救急車のサイレンの高い音は\qty{960}{Hz}です。救急車が時速\qty{60}{km}（$60\div3.6\approx\qty{16.7}{m/s}$）で走り、音の速さは\qty{340}{m/s}とします（中学で習う「約\qty{340}{m/s}」。気温\qty{15}{\celsius}くらいの値）。止まって聞く人には、
\[\text{近づくとき}\quad f'=\frac{340}{340-16.7}\times960\approx\qty{1010}{Hz}\]
\[\text{遠ざかるとき}\quad f'=\frac{340}{340+16.7}\times960\approx\qty{915}{Hz}\]
と聞こえます。目の前を通りすぎた瞬間に、\qty{1010}{Hz}から\qty{915}{Hz}へ、約1割も高さが下がります。音楽の音の高さでいうと、半音2つ分（全音）に近い下がり方です。近づくあいだは前の波面がつまって高く、遠ざかるあいだは後ろの波面がのびて低くなる。これが、通りすぎると低く聞こえるわけです。救急車といっしょに動いている運転手は、$v_{\mathrm o}=v_{\mathrm s}$ なので $f'=f$ で、いつも\qty{960}{Hz}を聞いています。

この本で使う\qty{20}{\celsius}の音の速さ\qty{343}{m/s}で計算すると、近づくとき\qty{1009}{Hz}、遠ざかるとき\qty{916}{Hz}です。違いは\qty{1}{Hz}以下で、聞き分けられません。音の速さの値の違い（約1\,\%）より、救急車の速さ（音の速さの約5\,\%）のほうが、ずっと大きく効くからです。

\begin{phmijika}[サイレンの大きさの決まり]
救急車のサイレンの大きさは、国の決まり（道路運送車両の保安基準の細目を定める告示 第231条）で、「自動車の前方\qty{20}{m}の位置において\qty{90}{dB}以上\qty{120}{dB}以下」と決められています。前方\qty{40}{m}では、距離が2倍なので約\qty{6}{dB}小さくなります（逆2乗の法則：波-7 ②）。サイレンは、\qty{960}{Hz}と\qty{770}{Hz}の音を交互に鳴らして「ピーポー」と聞こえるようにしてあり、どちらの音も、通りすぎるときにドップラー効果で約1割低くなります（\qty{770}{Hz}は、近づくとき約\qty{810}{Hz}、遠ざかるとき約\qty{734}{Hz}）。
\end{phmijika}

\phboxhead{速さが音よりずっと小さいとき}
$v_{\mathrm s}$ と $v_{\mathrm o}$ が $c$ よりずっと小さいときは、1次近似 $1/(1-x)\approx1+x$（5-17）を使うと、
\[f'=\Bigl(1-\frac{v_{\mathrm o}}{c}\Bigr)\frac{1}{1-v_{\mathrm s}/c}\,f\approx\Bigl(1-\frac{v_{\mathrm o}}{c}\Bigr)\Bigl(1+\frac{v_{\mathrm s}}{c}\Bigr)f\approx\Bigl(1+\frac{v_{\mathrm s}-v_{\mathrm o}}{c}\Bigr)f\]
となります（最後は、小さな数どうしのかけ算 $v_{\mathrm o}v_{\mathrm s}/c^2$ を無視した）。高さのずれの割合は、\textbf{近づく速さ÷音の速さ}です。救急車では $16.7\div340\approx0.049$ で、約5\,\% 高くなります（$960\times1.049\approx\qty{1007}{Hz}$。正しい式の\qty{1010}{Hz}とほぼ同じ）。

\begin{phbaisuru}{音源の速さ $v_{\mathrm s}$}
近づくときの高さのずれは、ほぼ2倍になります。時速\qty{60}{km}では\qty{960}{Hz}から約 $+\qty{50}{Hz}$、時速\qty{120}{km}では約 $+\qty{104}{Hz}$ です（正しい式で計算すると、2倍より少し大きい）。
\end{phbaisuru}

\phboxhead{風があるとき}
風が、音源から観測者の向きに速さ $w$ で吹いているときは、空気ごと動くので、地面から見た音の速さは $c+w$ になります。上の考えの $c$ を、すべて $c+w$ に置きかえて、
\pheq{f'=\frac{c+w-v_{\mathrm o}}{c+w-v_{\mathrm s}}\,f}{Hz}
\label{ph:波-8:kaze}です。音源も観測者も止まっていれば、風があっても $f'=f$ で、高さは変わりません（波長と速さが同じ割合で変わるだけ。例題11）。

\phboxhead{反射板があるとき}
止まっている音源から出た音を、近づいてくる物（速さ $v$）ではね返し、音源の所で聞きます。はね返す物は、まず「音源に近づく観測者」として $f_1=\dfrac{c+v}{c}f$ を受けとり、次に「観測者に近づく音源」として $f_1$ の音を出し直します。
\[f_2=\frac{c}{c-v}\,f_1=\frac{c+v}{c-v}\,f,\qquad f_2-f=\frac{2v}{c-v}\,f\approx\frac{2v}{c}\,f\]
はね返った音の高さのずれは、ほぼ $2vf/c$ で、1回の場合の2倍です（⑦の速さの計測のもと）。

\begin{phflow}[ドップラー効果]
\phflowrow{波面を数える（媒質に対して速さ $c$）}{一直線の上\\空気は止まっている}{$f'=\dfrac{c-v_{\mathrm o}}{c-v_{\mathrm s}}\,f$}
\phflowrow{$f'=\dfrac{c-v_{\mathrm o}}{c-v_{\mathrm s}}\,f$}{$v_{\mathrm s}$・$v_{\mathrm o}$ は $c$ より\\ずっと小さい}{$f'\approx\Bigl(1+\dfrac{v_{\mathrm s}-v_{\mathrm o}}{c}\Bigr)f$}
\phflowrow{$f'=\dfrac{c-v_{\mathrm o}}{c-v_{\mathrm s}}\,f$}{反射板で2回\\（受けて出す）}{$\Delta f\approx\dfrac{2v}{c}\,f$}
\end{phflow}

\phboxhead{斜めのドップラー効果}
救急車は、ふつう聞く人の真正面へ向かって走ってはきません。少し離れた道を、斜めの向きに通りすぎます。このとき効くのは、速さの\textbf{観測者へ向かう向きの成分}だけです。音を出したときの、救急車の進む向きと、救急車から観測者を見る向きのなす角を $\theta$ とすると、近づく速さの成分は $v_{\mathrm s}\cos\theta$（4-12 の内積：ある向きの成分）なので、
\pheq{f'=\frac{c}{c-v_{\mathrm s}\cos\theta}\,f}{Hz}
\label{ph:波-8:naname}です。真横（$\theta=90°$）では $\cos\theta=0$ で、高さは変わりません。

\begin{bkfigure}
\begin{tikzpicture}
\begin{axis}[bkplot, axis lines=left, width=106mm, height=48mm, xmin=-6, xmax=6.5, ymin=900, ymax=1022, xtick={-6,-4,-2,0,2,4,6}, ytick={915,960,1010}, xlabel={聞こえた時刻（s）}, ylabel={$f'$（Hz）}]
  \addplot[bkcurve, domain=-6:6, samples=240, variable=\u] ({\u+sqrt((16.667*\u)^2+100)/340}, {960*340/(340+16.667^2*\u/sqrt((16.667*\u)^2+100))});
  \addplot[bkOrange, line width=1.2pt, dash pattern=on 4pt off 2pt, domain=-6:6, samples=240, variable=\u] ({\u+sqrt((16.667*\u)^2+400)/340}, {960*340/(340+16.667^2*\u/sqrt((16.667*\u)^2+400))});
  \draw[bkGray, line width=0.5pt, dashed] (axis cs:-6,960) -- (axis cs:6.5,960);
  \node[font=\sffamily\footnotesize, anchor=south west, fill=white, inner sep=1pt] at (axis cs:1.2,985) {実線：離れ\qty{10}{m}};
  \node[font=\sffamily\footnotesize, anchor=north east, fill=white, inner sep=1pt] at (axis cs:-0.4,940) {破線：離れ\qty{20}{m}};
\end{axis}
\end{tikzpicture}
\bkfigcaption{道から\qty{10}{m}・\qty{20}{m}離れて聞く救急車の音（\qty{960}{Hz}、時速\qty{60}{km}、音の速さ\qty{340}{m/s}）。時刻0は、救急車がいちばん近くを通る時刻。近いほど、高さが急に下がる。}\label{ph:波-8:fig-naname}
\end{bkfigure}

\bksubsection{6}{衝撃波：音より速く動く}
\label{ph:波-8:sec6}
音源が音より速く動くと（$v_{\mathrm s}>c$）、音源は、自分が出した波面を追いこしてしまいます。前には波面ができず、たくさんの波面が、円すいの形の1枚の面に重なります。重なった所では圧力が急に高くなり、これを\textbf{衝撃波}といいます（熱-4 ⑨）。\bkidx{しょうげきは}{衝撃波}

\begin{bkfigure}
\begin{tikzpicture}[x=1mm, y=1mm]
  \begin{scope}
  \clip (-64,-26) rectangle (6,26);
  \foreach \n in {1,...,6} {\draw[bkBlue!70, line width=0.6pt] ({-2*4.5*\n},0) circle ({4.5*\n});}
  \draw[bkRed!80!black, line width=1.4pt] (0,0) -- (210:66);
  \draw[bkRed!80!black, line width=1.4pt] (0,0) -- (150:66);
  \end{scope}
  \fill[bkInk] (0,0) circle (1.3);
  \draw[bkInk, line width=0.8pt, -{Stealth[length=2.2mm]}] (2,0) -- (12,0);
  \node[font=\sffamily\footnotesize, anchor=west] at (12.5,0) {$v_{\mathrm s}$};
  \draw[bkInk, line width=0.5pt, dashed] (0,0) -- (-54,0);
  \draw[bkInk, line width=0.5pt] (-14,0) arc (180:150:14);
  \node[font=\sffamily\footnotesize, fill=white, inner sep=0.6pt] at (-18,3.4) {$\alpha$};
  \draw[bkGreen!60!black, line width=0.9pt] (-54,0) -- ($(-54,0)+(60:27)$);
  \node[font=\sffamily\footnotesize, anchor=north, fill=white, inner sep=0.6pt] at (-54,-1) {O};
  \node[anchor=west, font=\small, text width=48mm, align=left] at (18,4) {%
  O は、時間 $t$ 前に音源がいた所。\par\vspace{2pt}
  そのとき出た波面の半径は $ct$（O から斜め上へ引いた線）、音源が進んだ距離は $v_{\mathrm s}t$（破線）。\par\vspace{2pt}
  $\sin\alpha=\dfrac{ct}{v_{\mathrm s}t}=\dfrac{c}{v_{\mathrm s}}$};
\end{tikzpicture}
\bkfigcaption{音の2倍の速さで進む音源の波面と、円すいの形の衝撃波（太い2本の線の間の角が $2\alpha$）。}\label{ph:波-8:fig-mach}
\end{bkfigure}

\noindent 時間 $t$ 前に点Oで出た波面は、今は半径 $ct$ の球面です。そのあいだに、音源は O から $v_{\mathrm s}t$ 進んでいます。今の音源の位置から、すべての波面に接する線を引くと、それが衝撃波の面です。直角三角形で、
\pheq{\sin\alpha=\frac{c}{v_{\mathrm s}}=\frac{1}{\mathrm{Ma}}}{単位なし}
\label{ph:波-8:mach}です。$\alpha$ は円すいの半分の角、$\mathrm{Ma}=v_{\mathrm s}/c$ はマッハ数（熱-4 ⑦）です。\bkidx{まっはえんすい}{マッハ円すい}この円すいを\textbf{マッハ円すい}といいます。速いほど、円すいはとがります。$\mathrm{Ma}<1$ では $\sin\alpha>1$ となる角はなく、衝撃波はできません。

\begin{phmijika}[ソニックブーム・むちの音]
超音速の飛行機の円すいが地面を通りすぎると、地上の人には「ドーン」という大きな音（ソニックブーム）が聞こえます。飛行機が真上を通ったあと、少しおくれて聞こえます（例題10）。サーカスのむちの「パーン」という音は、むちの先が音より速く動くときにできる、小さな衝撃波だと考えられています。
\end{phmijika}

\begin{seikaku}
水面を進む船の後ろにも、V の字の波（引き波）ができます。ところが、深い水の上をふつうの速さで進む船の引き波の角（V の半分の角）は、船の速さによらず、約\qty{19.5}{\degree}になることが知られています（ケルビンの角。とても速いボートでは、もっと小さな角になることも観測されている）。深い水の波は、波長によって速さが違い、エネルギーが山の半分の速さで進む（波-7 ⑥）ので、マッハ円すいの式がそのままは使えないのです。
\end{seikaku}

\bksubsection{7}{ドップラーで速さをはかる}
\label{ph:波-8:sec7}
反射板の式 $\Delta f\approx2vf/c$ を逆に使うと、はね返った波の高さのずれ $\Delta f$ から、はね返した物の速さ $v$ がわかります。
\[v\approx\frac{c\,\Delta f}{2f}\]
\begin{itemize}[leftmargin=1.2\zw, itemsep=2pt, topsep=2pt, label={\color{bkTheme}●}]
\item \textbf{スピードガン}：電波（光と同じ速さ $c=3.00\times10^8$ m/s）を車やボールに当て、はね返った電波の振動数のずれから速さをはかる（例題8）。電波の速さはとても大きいので、ずれの割合はとても小さく、数百万分の1です。
\item \textbf{血流の検査}：超音波を血管に当てると、流れている血液の粒（赤血球）ではね返った音が、少しずれて戻ってきます。血管と超音波の向きが角 $\theta$ をなすときは、$\Delta f\approx2fv\cos\theta/c$（例題9）。ずれは数 kHz と、ちょうど耳に聞こえる高さなので、スピーカーから「ザッザッ」という音で血流を聞くこともできます。
\item \textbf{気象レーダー}：雨粒ではね返った電波のずれから、雨粒がレーダーに近づいているか遠ざかっているかがわかり、雲の中の風のようすがわかります。竜巻の兆しを見つけるのにも使われます。
\end{itemize}

\phboxhead{シミュレーション：波面を数えて確かめる}
⑤の式を、波面を1つずつ計算して確かめましょう。救急車が\qty{960}{Hz}の音を出しながら時速\qty{60}{km}で走ります。0秒から0.5秒までに出した波面について、0.5秒の時点での前と後ろの位置を計算し、となりどうしの間かくを平均します。
\begin{phcode}[Python]
c = 340.0              \# 音の速さ m/s
f = 960.0              \# サイレンの振動数 Hz
v = 60 / 3.6           \# 救急車の速さ m/s
t = 0.5                \# この時刻に、波面の位置を調べる s
T = 1 / f              \# 1つの波面を出す間かく s
mae = []               \# 前の波面の位置（m）
ato = []               \# 後ろの波面の位置（m）
n = 0
while n * T <= t:
    xs = v * n * T         \# 波面を出したときの救急車の位置
    R = c * (t - n * T)    \# その波面の、今の半径
    mae.append(xs + R)
    ato.append(xs - R)
    n = n + 1
kmae = (mae[0] - mae[-1]) / (n - 1)
kato = (ato[-1] - ato[0]) / (n - 1)
print("波面の数", n)
print("前", round(kmae, 4), "m", round(c / kmae, 1), "Hz")
print("後ろ", round(kato, 4), "m", round(c / kato, 1), "Hz")
\end{phcode}
\begin{phcodeout}
波面の数 481
前 0.3368 m 1009.5 Hz
後ろ 0.3715 m 915.1 Hz
\end{phcodeout}
前の間かく\qty{0.3368}{m}は $(c-v_{\mathrm s})/f=323.3\div960$、後ろの\qty{0.3715}{m}は $(c+v_{\mathrm s})/f$ と一致し、聞こえる高さ\qty{1009.5}{Hz}・\qty{915.1}{Hz}も⑤の式と同じです。$v$ を\qty{400}{m/s}（音より速い）にすると、前の位置の値が並ぶ順番が逆になり、前に波面がなくなること（⑥の衝撃波）も確かめられます。図\ref{ph:波-8:fig-doppler}・図\ref{ph:波-8:fig-mach}は、このように計算した円をかいたものです。

\begin{phmodel}[ドップラー効果の式]
\phignore{音源の大きさ（点とみなす）。空気の温度のむら。音が弱まること（高さには関係しない）。}
\phassume{空気（媒質）は地面に対して止まっているか、一定の風で動く。音源と観測者の速さは一定で、音より遅い。}
\phlimit{音源が観測者のすぐ近くを斜めに通るとき（$\theta$ が急に変わる：斜めの式を使う）。音より速いとき（衝撃波）。光や電波（媒質がない：相対性理論の式。ただし速さが小さければ同じ近似式になる）。}
\end{phmodel}

\begin{phmada}[宇宙の広がる速さは、いくら？]
遠くの銀河の光は、波長がのびて（赤い側へずれて）届きます。銀河が遠ざかっているからで、遠い銀河ほど速く遠ざかっています。これは、宇宙全体が広がっていることを示しています。その広がる速さを表す数（ハッブル定数）を、遠くの宇宙の光（宇宙が生まれて約38万年後の光）から求めると約67、近くの銀河の距離と光のずれから求めると約73（単位は km/s を、銀河までの距離の単位でわったもの）になり、測り方の誤差では説明しにくいほど食い違っています（2020年代の研究）。どちらかの測り方にまだ気づかれていないずれがあるのか、宇宙のしくみに新しいことがあるのかは、まだわかっていません。
\end{phmada}

\begin{phsaki}{光のドップラー効果と、波の重なり}
光や電波には、音にとっての空気のような媒質がありません。光の速さは、だれから見ても同じです（相対性理論）。そのため、光のドップラー効果の正しい式は、この章の式と少し違い、物理シリーズの『原子』（原-11）で学びます。速さが光よりずっと小さければ、この章の近似式 $\Delta f\approx vf/c$ と同じになります。光の干渉・回折（ヤングの実験、回折格子、円い穴の回折）は波-13 で、音のレベルと騒音は波-9 で、超音波の計測は波-10 で学びます。
\end{phsaki}

\begin{senmonsho}[音響・計測の専門書の書き方]
\begin{bkhonbun}
音源が速度 $v_{\mathrm s}$ で，観測者が速度 $v_{\mathrm o}$ で互いに近づくとき，観測される周波数は $f'=f(c+v_{\mathrm o})/(c-v_{\mathrm s})$ である．反射体が速度 $v$ で接近するときのドップラーシフトは $f_{\mathrm D}=2vf_0\cos\theta/c$ となる．
\end{bkhonbun}
\bkyaku{専門書は、「互いに近づく向き」をどちらも正にして書くことが多い（この章の $v_{\mathrm o}$ は、音源から観測者へ向かう向きが正なので、近づく観測者の $v_{\mathrm o}$ は負。どちらも同じ式）。ずれの $\Delta f$ を「ドップラーシフト」$f_{\mathrm D}$ とよぶ。$f_0$ は出した波の振動数。}
\end{senmonsho}

\begin{taisetsu}
\def\tsl#1{\makebox[7.6\zw][l]{\sffamily\bfseries\small\color{\bkTT}#1}}
\tsl{屈折の法則}$\dfrac{\sin\theta_1}{\sin\theta_2}=\dfrac{c_1}{c_2}=\dfrac{\lambda_1}{\lambda_2}$（振動数は変わらない）\\[3pt]
\tsl{うなり}1秒に $|f_1-f_2|$ 回\\[3pt]
\tsl{干渉}$|l_1-l_2|=m\lambda$ で強め合い、$(m+\tfrac12)\lambda$ で弱め合う\\[3pt]
\tsl{単スリット}最初の暗い所 $a\sin\theta=\lambda$\\[3pt]
\tsl{ドップラー}$f'=\dfrac{c-v_{\mathrm o}}{c-v_{\mathrm s}}f$、斜めは $v\cos\theta$、反射板は $\Delta f\approx\dfrac{2v}{c}f$\\[3pt]
\tsl{衝撃波}$\sin\alpha=\dfrac{c}{v_{\mathrm s}}=\dfrac{1}{\mathrm{Ma}}$
\end{taisetsu}

\begin{kotoba}
\small
\bkword{ホイヘンスの原理}{波面の各点が素元波の源になり、その包絡面が次の波面になる、という考え。}
\bkword{うなり}{振動数の近い2つの音が重なって、音の大きさが1秒に $|f_1-f_2|$ 回ゆれること。}
\bkword{干渉}{波が重なって、強め合ったり弱め合ったりすること。}
\bkword{回折}{波が、すきまを通ったり物のはしを通ったりして、かげへ回りこむこと。}
\bkword{ドップラー効果}{音源や観測者が動くと、聞こえる振動数が変わること。}
\bkword{衝撃波}{音より速く動く物の前にできる、圧力が急に変わる面。}
\end{kotoba}

\bkkigocards{\bkkigocard{$v_{\mathrm s}$・$v_{\mathrm o}$}{ブイエス・ブイオー}{音源（source）と観測者（observer）の速さ。この巻では、音源から観測者へ向かう向きを正とする。音の速さは $c$（高校の教科書は $V$）}{velocity of source / observer}{波-8}}
{\bkkigocard{$a$}{エー}{スリット（すきま）の幅。2つの波源・スリットの間かくは $d$。干渉の次数は $m$}{slit width}{波-8}}

%% ======================== 例題 ========================
\Needspace{10\baselineskip}\bksection{reidai}{例題でたしかめよう}
\label{ph:波-8:ex1}
\begin{phreidai}[救急車のサイレン（はてな？）]
\qty{960}{Hz}のサイレンを鳴らす救急車が、時速\qty{60}{km}でまっすぐ近づいてきて、目の前を通りすぎ、遠ざかっていきます。止まって聞く人に聞こえる振動数を、近づくときと遠ざかるときについて求めましょう。音の速さは\qty{340}{m/s}とします。また、前と後ろの波長も求めましょう。
\phbamen
音源が動き、観測者が止まっている、ドップラー効果の場面です。
\phmushi
風はないとします。救急車は観測者の真正面へ近づき、真後ろへ遠ざかる（一直線の上）とします（斜めのときは例題7）。
\phzu
図\ref{ph:波-8:fig-doppler}。$f=\qty{960}{Hz}$、$c=\qty{340}{m/s}$、$v_{\mathrm s}=60\div3.6\approx\qty{16.7}{m/s}$、$v_{\mathrm o}=0$。
\phriyu
音源が動くと、波面の間かくが変わります。観測者は止まっているので、$f'=c/\lambda'$（⑤）。
\phtateru
\[\lambda'=\frac{c-v_{\mathrm s}}{f},\qquad f'=\frac{c}{c-v_{\mathrm s}}\,f\]
遠ざかるときは、$v_{\mathrm s}$ を $-v_{\mathrm s}$ にします。
\phtoku
近づくとき：\[\lambda'=\frac{340-16.7}{960}\approx\qty{0.337}{m},\qquad f'=\frac{340}{323.3}\times960\approx\qty{1010}{Hz}\]
遠ざかるとき：\[\lambda'=\frac{340+16.7}{960}\approx\qty{0.372}{m},\qquad f'=\frac{340}{356.7}\times960\approx\qty{915}{Hz}\]
\phtashikame
1次近似（⑤）では $960\times(1\pm0.049)$ で、\qty{1007}{Hz}と\qty{913}{Hz}。正しい式とほぼ合います。\\
極端な場合：$v_{\mathrm s}=0$ なら $f'=f$。止まっている救急車のサイレンは\qty{960}{Hz}に聞こえます。\\
止まっているときの波長は $340\div960\approx\qty{0.354}{m}$ で、前ではそれより短く、後ろでは長くなっています。
\kotae 近づくとき約\qty{1010}{Hz}（波長\qty{0.337}{m}）、遠ざかるとき約\qty{915}{Hz}（波長\qty{0.372}{m}）
\end{phreidai}
\begin{jibun}
音の速さを\qty{340}{m/s}とします。（1）例題1の救急車が時速\qty{100}{km}で走るとき、近づくときと遠ざかるときの振動数を求めましょう。（2）時速\qty{80}{km}の救急車の低い音（\qty{770}{Hz}）ではどうですか。
\bkjans{（1）約\qty{1045}{Hz}と約\qty{887}{Hz}。（2）約\qty{824}{Hz}と約\qty{723}{Hz}}
\end{jibun}

\label{ph:波-8:ex2}
\begin{phreidai}[チューニングのうなり]
おんさの\qty{440}{Hz}の音と、ギターの弦の\qty{442}{Hz}の音を同時に鳴らします。1秒に何回うなりますか。聞こえる音の高さはいくらですか。
\phbamen
振動数の近い2つの音が重なる場面です。
\phmushi
2つの音の振幅は同じで、ほかの高さの音（倍音）はないとします。
\phzu
$f_1=\qty{442}{Hz}$、$f_2=\qty{440}{Hz}$。
\phriyu
和を積に直すと、まん中の高さの音の振幅が、ゆっくり変わります。音が大きくなる回数は $|f_1-f_2|$（②）。
\phtateru
\[\text{うなりの回数}=|f_1-f_2|,\qquad\text{聞こえる高さ}=\frac{f_1+f_2}{2}\]
\phtoku
\[|442-440|=2\ \text{（1秒に2回）},\qquad\frac{442+440}{2}=\qty{441}{Hz}\]
\phtashikame
4-9⑦で計算した結果と同じです。うなりの周期（大きくなってから次に大きくなるまで）は $1\div2=\qty{0.5}{s}$。\\
極端な場合：弦を\qty{440}{Hz}に合わせると、差が0になり、うなりは消えます。
\kotae 1秒に2回、\qty{441}{Hz}の音に聞こえる
\end{phreidai}
\begin{jibun}
（1）\qty{110}{Hz}と\qty{111.5}{Hz}の音を重ねると、1秒に何回うなりますか。うなりの周期は何秒ですか。（2）\qty{440}{Hz}のおんさと弦を鳴らすと、1秒に3回うなりました。弦を少し強く張る（高くする）と、うなりがゆっくりになりました。はじめの弦の振動数はいくらでしたか。
\bkjans{（1）1秒に1.5回、周期 約\qty{0.67}{s}。（2）高くしてうなりが減ったので、はじめはおんさより低い。$440-3=\qty{437}{Hz}$}
\end{jibun}

\label{ph:波-8:ex3}
\begin{phreidai}[2台のスピーカー]
2台のスピーカーから、同じ位相で\qty{686}{Hz}の音が出ています。点Pは、一方のスピーカーから\qty{3.00}{m}、もう一方から\qty{3.25}{m}の所にあります。点Pでは、音は強め合いますか、弱め合いますか。音の速さは\qty{343}{m/s}とします。
\phbamen
2つの波源からの音が重なる、干渉の場面です。
\phmushi
道のりによる音の弱まりの違い（$3.00$ と $3.25$ はほぼ同じ）と、部屋の壁ではね返る音を無視します。
\phzu
図\ref{ph:波-8:fig-kanshou}。$l_1=\qty{3.00}{m}$、$l_2=\qty{3.25}{m}$。
\phriyu
強め合うか弱め合うかは、道のりの差が波長の何倍かで決まります（③）。
\phtateru
\[\lambda=\frac{c}{f},\qquad\frac{|l_1-l_2|}{\lambda}\ \text{が整数なら強め合い、整数}+\tfrac12\ \text{なら弱め合い}\]
\phtoku
\[\lambda=\frac{343}{686}=\qty{0.500}{m},\qquad\frac{|3.00-3.25|}{0.500}=\frac{0.25}{0.500}=0.5\]
道のりの差が半波長なので、弱め合います。
\phtashikame
単位：m÷m で単位なし。\\
2台を\qty{1.0}{m}はなし、\qty{3}{m}前で横に歩くと、まん中（差0）から約\qty{0.78}{m}横で、差がちょうど\qty{0.25}{m}になります。やってみよう②の「小さく聞こえる所」です。
\kotae 道のりの差は半波長なので、弱め合う
\end{phreidai}
\begin{jibun}
例題3の点Pで、音の振動数を（1）\qty{343}{Hz}、（2）\qty{1372}{Hz}にすると、強め合いますか、弱め合いますか。
\bkjans{（1）$\lambda=\qty{1.00}{m}$、差は0.25波長で、強め合いと弱め合いの中間（$2A\cos45°\approx1.4A$）。（2）$\lambda=\qty{0.25}{m}$、差はちょうど1波長で、強め合う}
\end{jibun}

\label{ph:波-8:ex4}
\begin{phreidai}[海岸の波は、なぜ岸に平行になる？]
沖から、波面が岸の線と\qty{30}{\degree}の角をなして（入射角\qty{30}{\degree}で）波が来ます。深さ\qty{4.0}{m}の所から、深さ\qty{1.0}{m}の所へ入るとき、屈折角を求めましょう。浅い水の波の速さは $c=\sqrt{gh}$（$h$ は深さ。波-5 ⑤）、$g=\qty{9.8}{m/s^2}$ とします。
\phbamen
深さが変わる所を、波が斜めに進む場面です。
\phmushi
深さは、境目で急に変わるとします（実際はなだらかに変わる）。波は小さく、風はないとします。
\phzu
図\ref{ph:波-8:fig-kussetsu}。$h_1=\qty{4.0}{m}$、$h_2=\qty{1.0}{m}$、$\theta_1=\qty{30}{\degree}$。
\phriyu
速さの違う所へ斜めに入るので、屈折の法則（①）。速さは深さで決まります。
\phtateru
\[c=\sqrt{gh},\qquad\frac{\sin\theta_1}{\sin\theta_2}=\frac{c_1}{c_2}\]
\phtoku
\[c_1=\sqrt{9.8\times4.0}\approx\qty{6.26}{m/s},\qquad c_2=\sqrt{9.8\times1.0}\approx\qty{3.13}{m/s}\]
\[\sin\theta_2=\sin30°\times\frac{3.13}{6.26}=0.5\times0.5=0.25,\qquad\theta_2\approx\qty{14}{\degree}\]
\phtashikame
深さが $\tfrac14$ なので速さは半分、$\sin$ も半分。浅い所に入るほど垂線（岸に直角な線）に近づき、波面は岸に平行に近づきます。沖ではいろいろな向きから来る波も、岸の近くでは、ほとんど岸に平行に打ち寄せるのはこのためです。
\kotae 約\qty{14}{\degree}（波面が岸に平行に近づく）
\end{phreidai}
\begin{jibun}
深さ\qty{8.0}{m}の所から、入射角\qty{45}{\degree}で、深さ\qty{0.50}{m}の所へ入る波の屈折角を求めましょう。
\bkjans{$c_2/c_1=\sqrt{0.50\div8.0}=0.25$、$\sin\theta_2=0.707\times0.25\approx0.18$、約\qty{10}{\degree}}
\end{jibun}

\label{ph:波-8:ex5}
\begin{phreidai}[声は回りこみ、光は回りこまない]
（1）\qty{500}{Hz}の声が、幅\qty{0.80}{m}の戸口を通るとき、最初の暗い（静かな）向きの角を求めましょう。音の速さは\qty{343}{m/s}です。（2）波長\qty{633}{nm}（$633\times10^{-9}$\,m）の赤いレーザーの光が、幅\qty{0.10}{mm}のすきまを通るときはどうですか。
\phbamen
すきまを通った波の広がり方（回折）を、音と光で比べる場面です。
\phmushi
すきまの両側の壁は、波を通さないとします。すきまに来る波は平面波とします。
\phzu
図\ref{ph:波-8:fig-slit}。（1）$\lambda=343\div500$、$a=\qty{0.80}{m}$。（2）$\lambda=633\times10^{-9}$\,m、$a=1.0\times10^{-4}$\,m。
\phriyu
最初の暗い所は $a\sin\theta=\lambda$（④）。
\phtateru
\[\sin\theta=\frac{\lambda}{a}\]
\phtoku
（1）\[\lambda=\frac{343}{500}\approx\qty{0.686}{m},\qquad\sin\theta=\frac{0.686}{0.80}\approx0.86,\qquad\theta\approx\qty{59}{\degree}\]
（2）\[\sin\theta=\frac{633\times10^{-9}}{1.0\times10^{-4}}=6.33\times10^{-3},\qquad\theta\approx\qty{0.36}{\degree}\]
\phtashikame
声は\qty{59}{\degree}まで大きく広がり、戸口の横にいても聞こえます。光は、すきまをせまくしても\qty{0.36}{\degree}しか広がりません。\qty{2}{m}先のかべでは、まん中から $2\times6.33\times10^{-3}\approx\qty{0.013}{m}$（約\qty{13}{mm}）の所が最初の暗い所です。戸口（\qty{0.80}{m}）を通る光なら、$\sin\theta\approx7\times10^{-7}$ で、回りこみは見えません。
\kotae （1）約\qty{59}{\degree}、（2）約\qty{0.36}{\degree}
\end{phreidai}
\begin{jibun}
（1）例題5（2）のすきまの幅を\qty{0.050}{mm}にすると、最初の暗い所の角と、\qty{2}{m}先のかべでの位置はどうなりますか。（2）\qty{3000}{Hz}の音が幅\qty{0.80}{m}の戸口を通るとき、最初の静かな向きの角を求めましょう（音の速さ\qty{343}{m/s}）。
\bkjans{（1）約\qty{0.73}{\degree}、まん中から約\qty{25}{mm}（2倍に広がる）。（2）$\lambda\approx\qty{0.114}{m}$、$\sin\theta\approx0.14$、約\qty{8}{\degree}}
\end{jibun}

\label{ph:波-8:ex6}
\begin{phreidai}[電車の中で聞く、踏切の音]
時速\qty{72}{km}で走る電車に乗っています。踏切の警報音（ここでは\qty{700}{Hz}とする）は、近づくときと遠ざかるときに、何 Hz に聞こえますか。音の速さは\qty{343}{m/s}とします。
\phbamen
音源が止まっていて、観測者が動く場面です。
\phmushi
風はないとします。電車は踏切に向かってまっすぐ進む（一直線）とします。
\phzu
$f=\qty{700}{Hz}$、$v_{\mathrm s}=0$、電車の速さ $72\div3.6=\qty{20}{m/s}$。近づくときは、音源から観測者へ向かう向きと反対に動くので $v_{\mathrm o}=\qty{-20}{m/s}$、遠ざかるときは $v_{\mathrm o}=\qty{+20}{m/s}$。
\phriyu
波面の間かくは変わらず、波面とすれ違う速さが変わります（⑤）。
\phtateru
\[f'=\frac{c-v_{\mathrm o}}{c}\,f\]
\phtoku
\[\text{近づくとき}\quad f'=\frac{343+20}{343}\times700\approx\qty{741}{Hz}\]
\[\text{遠ざかるとき}\quad f'=\frac{343-20}{343}\times700\approx\qty{659}{Hz}\]
\phtashikame
ずれは $\pm700\times20\div343\approx\pm\qty{41}{Hz}$ で、近づくのも遠ざかるのも同じ大きさです（観測者が動くときは、正しい式でも左右対称）。音源が動くとき（例題1）は、近づくほうのずれが少し大きくなります。
\kotae 近づくとき約\qty{741}{Hz}、遠ざかるとき約\qty{659}{Hz}
\end{phreidai}
\begin{jibun}
（1）例題6の電車が時速\qty{90}{km}のときはどうですか。（2）救急車（\qty{960}{Hz}、時速\qty{60}{km}）と、時速\qty{40}{km}の車が、互いに近づいています。車の中の人に聞こえる振動数を求めましょう（音の速さ\qty{340}{m/s}）。
\bkjans{（1）約\qty{751}{Hz}と約\qty{649}{Hz}。（2）$v_{\mathrm s}=\qty{16.7}{m/s}$、$v_{\mathrm o}=\qty{-11.1}{m/s}$ で、約\qty{1042}{Hz}}
\end{jibun}

\label{ph:波-8:ex7}
\begin{phreidai}[斜めに通りすぎる救急車]
救急車（\qty{960}{Hz}、時速\qty{60}{km}）が、道から\qty{10}{m}離れて立つ人の前を通りすぎます。救急車の進む向きと、救急車から人を見る向きのなす角が\qty{60}{\degree}のとき（近づいている）に出た音は、何 Hz に聞こえますか。音の速さは\qty{340}{m/s}とします。また、\qty{960}{Hz}がそのまま聞こえるのは、いつ出た音ですか。
\phbamen
音源が、観測者の正面でなく、斜めに近づく場面です。
\phmushi
風はないとします。救急車は点とみなします。
\phzu
図\ref{ph:波-8:fig-naname}。$\theta=\qty{60}{\degree}$、$v_{\mathrm s}=\qty{16.7}{m/s}$。
\phriyu
観測者へ向かう速さの成分 $v_{\mathrm s}\cos\theta$ だけが効きます（⑤の斜めの式）。
\phtateru
\[f'=\frac{c}{c-v_{\mathrm s}\cos\theta}\,f\]
\phtoku
\[f'=\frac{340}{340-16.7\times0.5}\times960=\frac{340}{331.7}\times960\approx\qty{984}{Hz}\]
$\cos\theta=0$（$\theta=\qty{90}{\degree}$）のとき $f'=f$ です。これは、救急車がちょうど人の真横（いちばん近い所）にいるときに出た音で、$10\div340\approx\qty{0.03}{s}$ 後に人に届きます。
\phtashikame
\qty{984}{Hz}は、正面から近づくとき（\qty{1010}{Hz}）と止まっているとき（\qty{960}{Hz}）のあいだです。図\ref{ph:波-8:fig-naname}のように、高さは、通りすぎる前後の短い時間に、なめらかに下がります。
\kotae 約\qty{984}{Hz}。\qty{960}{Hz}は、救急車が真横を通るときに出た音（約\qty{0.03}{s}後に聞こえる）
\end{phreidai}
\begin{jibun}
例題7で、救急車が通りすぎたあと、角が\qty{120}{\degree}のときに出た音は何 Hz に聞こえますか。また、人が道から\qty{20}{m}離れていると、\qty{960}{Hz}が聞こえるのは、救急車が真横を通ってから何秒後ですか。
\bkjans{$\cos120°=-0.5$ で、$340\div(340+8.3)\times960\approx\qty{937}{Hz}$。$20\div340\approx\qty{0.06}{s}$ 後}
\end{jibun}

\label{ph:波-8:ex8}
\begin{phreidai}[スピードガン]
スピードガンから、24 GHz 帯の電波（ここでは振動数 $24.15\times10^9$\,Hz とする）を出し、\qty{40}{m/s}（時速\qty{144}{km}）で近づいてくる野球のボールではね返らせます。はね返った電波の振動数のずれ $\Delta f$ を求めましょう。電波の速さは $3.00\times10^8$\,m/s です。
\phbamen
動く物で波をはね返し、高さのずれから速さをはかる場面です。
\phmushi
ボールはまっすぐスピードガンに向かってくるとします。光の速さに近い速さではないので、近似の式を使います。
\phzu
$f=24.15\times10^9$\,Hz、$v=\qty{40}{m/s}$、$c=3.00\times10^8$\,m/s。
\phriyu
反射板が近づくとき、ずれは $\Delta f\approx2vf/c$（⑤の反射板の式）。
\phtateru
\[\Delta f\approx\frac{2v}{c}\,f\]
\phtoku
\[\Delta f\approx\frac{2\times40}{3.00\times10^8}\times24.15\times10^9=\frac{80\times24.15\times10^9}{3.00\times10^8}\approx6.4\times10^3\ \text{Hz}\]
\phtashikame
ずれの割合 $2v/c\approx2.7\times10^{-7}$ は、とても小さい値です。スピードガンは、はね返った電波ともとの電波を重ねて、そのうなり（②）の回数として $\Delta f$ を数えます。\\
極端な場合：$v$ が2倍なら $\Delta f$ も2倍です。止まったボールでは $\Delta f=0$。
\kotae 約\qty{6.4}{kHz}
\end{phreidai}
\begin{jibun}
例題8のスピードガンで、ずれが\qty{4.0}{kHz}でした。ボールの速さは何 m/s、時速何 km ですか。
\bkjans{$v\approx c\,\Delta f/(2f)$ より、約\qty{25}{m/s}（時速約\qty{89}{km}）}
\end{jibun}

\label{ph:波-8:ex9}
\begin{phreidai}[血液の流れの速さ]
振動数\qty{5.0}{MHz}の超音波を、血管に\qty{60}{\degree}の角で当てます。血液が\qty{0.50}{m/s}で流れているとき、はね返った音のずれ $\Delta f$ を求めましょう。体の中の音の速さは\qty{1540}{m/s}とします。
\phbamen
体の中を流れる血液の速さを、外から超音波ではかる場面です。
\phmushi
血液の粒はみな同じ速さで流れているとし、体の中での音の吸収は高さに関係しないので無視します。
\phzu
$f=5.0\times10^6$\,Hz、$v=\qty{0.50}{m/s}$、$\theta=\qty{60}{\degree}$、$c=\qty{1540}{m/s}$。
\phriyu
反射板の式に、斜めの成分 $v\cos\theta$ を入れます（⑤・⑦）。
\phtateru
\[\Delta f\approx\frac{2fv\cos\theta}{c}\]
\phtoku
\[\Delta f\approx\frac{2\times5.0\times10^6\times0.50\times0.5}{1540}\approx\qty{1.6e3}{Hz}\]
\phtashikame
\qty{1.6}{kHz}は、耳で聞ける高さです。\\
極端な場合：超音波を血管に直角（$\theta=\qty{90}{\degree}$）に当てると、$\cos\theta=0$ でずれが出ず、速さがはかれません。検査では、斜めに当てるのが大事です。
\kotae 約\qty{1.6}{kHz}
\end{phreidai}
\begin{jibun}
振動数\qty{3.0}{MHz}の超音波を\qty{45}{\degree}で当てたら、ずれが\qty{2.6}{kHz}でした。血液の速さを求めましょう（音の速さ\qty{1540}{m/s}）。
\bkjans{$v\approx c\,\Delta f/(2f\cos\theta)$ より、約\qty{0.94}{m/s}}
\end{jibun}

\label{ph:波-8:ex10}
\begin{phreidai}[マッハ2の飛行機の円すい]
マッハ2で飛ぶ飛行機のマッハ円すいの、半分の角 $\alpha$ を求めましょう。また、この飛行機が高さ\qty{10}{km}をまっすぐ水平に飛んでいるとき、真下の人にソニックブームが聞こえるのは、飛行機が真上を通ってから何秒後ですか。高い空の音の速さは\qty{295}{m/s}とします（熱-4 例題9）。
\phbamen
音より速い物が作る衝撃波の場面です。
\phmushi
空気の温度は高さによらず一定（音の速さが一定）とし、衝撃波の面はまっすぐな円すいとします。
\phzu
図\ref{ph:波-8:fig-mach}。$\mathrm{Ma}=2$、高さ $h=\qty{10000}{m}$、$c=\qty{295}{m/s}$。
\phriyu
円すいの角は $\sin\alpha=1/\mathrm{Ma}$（⑥）。円すいの面が人に届くのは、飛行機が人より $h/\tan\alpha$ 先へ進んだときです。
\phtateru
\[\sin\alpha=\frac{1}{\mathrm{Ma}},\qquad\text{先へ進む距離}=\frac{h}{\tan\alpha},\qquad\text{時間}=\frac{h/\tan\alpha}{\mathrm{Ma}\,c}\]
\phtoku
\[\sin\alpha=\frac12,\qquad\alpha=\qty{30}{\degree}\]
\[\frac{10000}{\tan30°}\approx\qty{17300}{m},\qquad\frac{17300}{2\times295}\approx\qty{29}{s}\]
\phtashikame
単位：m÷(m/s)＝s。\\
大きさ：真上を通ってから30秒ほどして、もう飛行機が\qty{17}{km}も先にいるころに「ドーン」と聞こえます。\\
極端な場合：Ma が1に近いと $\alpha$ は\qty{90}{\degree}に近く、円すいは平らな面になります。
\kotae $\alpha=\qty{30}{\degree}$、約\qty{29}{s}後
\end{phreidai}
\begin{jibun}
（1）マッハ円すいの半分の角が\qty{45}{\degree}のとき、マッハ数はいくらですか。（2）マッハ3では、半分の角は何度ですか。
\bkjans{（1）$1\div\sin45°\approx1.4$。（2）$\sin\alpha=\tfrac13$、約\qty{19}{\degree}}
\end{jibun}

\label{ph:波-8:ex11}
\begin{phreidai}[風が吹いているとき]
音の速さは\qty{340}{m/s}とします。（1）止まっている救急車（\qty{960}{Hz}）から、止まっている人の向きへ、風が\qty{10}{m/s}で吹いています。人に聞こえる振動数と、救急車の前（風下）と後ろ（風上）の波長を求めましょう。（2）この風の中を、救急車が人に向かって時速\qty{60}{km}で近づくと、何 Hz に聞こえますか。
\phbamen
媒質（空気）そのものが動いている場面です。
\phmushi
風の速さは、どこでも同じで一定とします。
\phzu
$w=\qty{10}{m/s}$（音源から観測者の向きが正）、$c=\qty{340}{m/s}$。
\phriyu
風があるときは、地面から見た音の速さが $c+w$（風下）・$c-w$（風上）になります。⑤の式の $c$ を $c+w$ に置きかえます。
\phtateru
\[f'=\frac{c+w-v_{\mathrm o}}{c+w-v_{\mathrm s}}\,f\]
\phtoku
（1）$v_{\mathrm s}=v_{\mathrm o}=0$ なので $f'=\dfrac{350}{350}\times960=\qty{960}{Hz}$。\\
風下の波長 $350\div960\approx\qty{0.365}{m}$、風上の波長 $330\div960\approx\qty{0.344}{m}$。\\
（2）\[f'=\frac{350}{350-16.7}\times960\approx\qty{1008}{Hz}\]
\phtashikame
（1）では、波長は変わっても、波の速さも同じ割合で変わるので、1秒に届く波面の数は変わりません。（2）は、風がないとき（\qty{1010}{Hz}）より\qty{2}{Hz}ほど低くなるだけです。風は、高さのずれをほんの少し変えるだけです。
\kotae （1）\qty{960}{Hz}（変わらない）、波長は風下\qty{0.365}{m}・風上\qty{0.344}{m}。（2）約\qty{1008}{Hz}
\end{phreidai}
\begin{jibun}
例題11（2）で、風が反対に、人から救急車の向きへ\qty{10}{m/s}で吹いていると、何 Hz に聞こえますか。
\bkjans{$c+w=\qty{330}{m/s}$ として、$330\div(330-16.7)\times960\approx\qty{1011}{Hz}$}
\end{jibun}

\label{ph:波-8:ex12}
\begin{phreidai}[コウモリが聞くこだま]
コウモリが\qty{40}{kHz}の超音波（日本の町でよく見るアブラコウモリが獲物を探すときの声くらい）を出しながら、\qty{5.0}{m/s}で止まっている壁へまっすぐ飛んでいきます。コウモリが聞くこだまの振動数と、出した音とのずれを求めましょう。音の速さは\qty{343}{m/s}とします。
\phbamen
音源と観測者が同じ物（コウモリ）で、はね返す物が止まっている場面です。
\phmushi
風はないとし、コウモリは壁に向かってまっすぐ飛ぶとします。
\phzu
$f=\qty{40e3}{Hz}$、$v=\qty{5.0}{m/s}$、$c=\qty{343}{m/s}$。
\phriyu
壁は、近づく音源からの音を受けとり（$f_1=\tfrac{c}{c-v}f$）、そのまま出し直します。それを、壁に近づく観測者（コウモリ）が聞きます（$f_2=\tfrac{c+v}{c}f_1$）。2回のドップラー効果で、⑤の反射板と同じ形になります。
\phtateru
\[f_2=\frac{c+v}{c-v}\,f\]
\phtoku
\[f_2=\frac{348}{338}\times40\times10^3\approx\qty{41.18e3}{Hz},\qquad f_2-f\approx\qty{1.18e3}{Hz}\]
\phtashikame
近似の式 $2vf/c=2\times5.0\times\num{40000}\div343\approx\qty{1.17e3}{Hz}$ とほぼ合います。コウモリは、こだまの高さのずれ（約\qty{1.2}{kHz}）から、どれくらいの速さで近づいているかを知ることができます（波-10）。
\kotae 約\qty{41.2}{kHz}、ずれは約\qty{1.2}{kHz}
\end{phreidai}
\begin{jibun}
例題12のコウモリ（\qty{40}{kHz}、\qty{5.0}{m/s}）が、\qty{1.0}{m/s}で遠ざかるガを追いかけています。こだまのずれは、近似の式 $\Delta f\approx2(v_{\text{近づく速さ}})f/c$ でいくらですか。近づく速さは、コウモリとガの速さの差とします。
\bkjans{近づく速さは $5.0-1.0=\qty{4.0}{m/s}$ なので、約\qty{0.93}{kHz}}
\end{jibun}

\begin{phmijika}[ずれを打ち消すコウモリ]
キクガシラコウモリのなかまは、自分が飛ぶ速さで生じるこだまのずれを打ち消すように、出す音の高さを少し下げ、こだまがいつも同じ高さ（耳がいちばんよく聞こえる高さ）で返ってくるようにしていることが知られています。ドップラー効果を、体の中で計算しているかのようです。
\end{phmijika}

%% ======================== まちがえやすいところ ========================
\Needspace{11\baselineskip}\bksection{tsumazuki}{まちがえやすいところをチェック}
\begin{tsumazuki}{うなりの回数は、差の半分？}
\bkng{振幅が $2A\cos\bigl(2\pi\tfrac{f_1-f_2}{2}t\bigr)$ なので、うなりは1秒に $\tfrac{f_1-f_2}{2}$ 回。}
\bkok{音の大きさは振幅の\textbf{大きさ}で決まり、$\cos$ の1周期に2回大きくなるので、1秒に $|f_1-f_2|$ 回。}
4-9⑦のグラフで、大きくなる回数を数えて確かめましょう（②）。
\end{tsumazuki}

\begin{tsumazuki}{近づくと高く聞こえるのは、音が大きくなるから？}
\bkng{救急車が近づくと音が大きくなるので、高く聞こえる。}
\bkok{高さは、1秒に耳に届く波面の数で決まる。近づくと波面の間かくがつまるので、高くなる。大きさとは別。}
真横を通る瞬間、音はいちばん大きいのに、高さは\qty{960}{Hz}のままです（⑤・例題7）。
\end{tsumazuki}

\begin{tsumazuki}{音源が動くのと、聞く人が動くのは同じ？}
\bkng{近づく速さが同じなら、どちらが動いても高さは同じ。}
\bkok{空気に対して動くのがどちらかで、少し違う。救急車が\qty{16.7}{m/s}で近づくと約\qty{1010}{Hz}、人が\qty{16.7}{m/s}で近づくと約\qty{1007}{Hz}。}
音は空気に対して速さ $c$ で進むので、空気に対する動きが大事です。速さが音よりずっと小さければ、ほぼ同じになります（⑤の近似式）。
\end{tsumazuki}

\begin{tsumazuki}{風が吹くと、高さが変わる？}
\bkng{風下にいると、音が速く届くので、高く聞こえる。}
\bkok{音源も人も止まっていれば、風があっても高さは変わらない。速さも波長も同じ割合で変わるから。}
音源や人が動くときだけ、風が少し効きます（例題11）。
\end{tsumazuki}

\begin{tsumazuki}{光は回折しない？}
\bkng{光はまっすぐ進むので、回りこまない。}
\bkok{光も波なので回折する。波長がとても短いので、ふつうのすきまでは回りこみが小さく、目立たないだけ。}
すきまを\qty{0.1}{mm}くらいにすると、光も広がるのが見えます（例題5）。
\end{tsumazuki}

%% ======================== 工学ではこう使う ========================
\Needspace{18\baselineskip}\bksection{phgear}{工学ではこう使う}
\begin{phkougaku}[重なりと高さのずれを使う]
\textbf{ドップラーレーダー}：気象庁や空港のドップラーレーダーは、雨粒ではね返った電波の高さのずれから、雲の中の風の速さと向きを調べます。飛行機の離着陸を危なくする、急に吹き下ろす風（ダウンバースト）を見つけるのにも使われます。

\textbf{超音波のドップラー計}：病院の血流の検査、工場の管の中を流れる液体の速さをはかる流量計など、「はね返った波の高さのずれ」から速さを読む装置は、たくさんあります。

\textbf{騒音を減らす干渉}：ノイズキャンセリングのヘッドホンや、車の中のこもった音を消す装置は、逆の位相の音を重ねて弱め合わせます。ただし、打ち消せるのは、波長が長く、音の場所による違いが小さい、低い音が中心です。

\begin{phketa}[波の重なりと高さのずれの数の目安]
\phketaline{音の波長}{\qty{20}{Hz}で\qty{17}{m}、\qty{1}{kHz}で\qty{34}{cm}、\qty{20}{kHz}で\qty{1.7}{cm}}
\phketaline{光の波長}{約 \qtyrange{380}{780}{nm}（紫〜赤。$0.38\times10^{-6}$〜$0.78\times10^{-6}$\,m）}
\phketaline{ずれの割合 $v/c$}{救急車と音：約5\,\%、ボールと電波：約 $10^{-7}$}
\phboxhead{答えが桁で合っているか}
ドップラー効果で、音の高さのずれが2倍をこえるような答えが出たら、速さが音の速さに近いか、式の符号をまちがえています。車や人の速さでは、音の高さのずれは数\,\% から十数\,\% です。回折の角が\qty{90}{\degree}をこえる（$\sin\theta>1$）ときは、波がすべての向きに広がる、という意味です。
\end{phketa}
\end{phkougaku}
\phstation{音響工学,レーダー・リモートセンシング,医用超音波,航空工学}{干渉、回折、ドップラーシフト、マッハ数、ホイヘンスの原理、波面}

\begin{column}{ドップラーと、列車の上のラッパ}
オーストリアのクリスチャン・ドップラーは、1842年、「動く星の光の色は、星の動きで変わるはずだ」という論文を書きました。この考えを音で確かめたのが、オランダのボイス・バロットです。1845年、走る列車の上でラッパ吹きに一定の高さの音を出してもらい、線路のそばで、音の高さのよくわかる音楽家たちに聞いてもらいました。列車が近づくときは高く、遠ざかるときは低く聞こえ、ドップラーの考えが正しいことがわかりました。ドップラーが考えた光のずれは、そのあと、星や銀河の動きを調べる大事な道具になりました。
\end{column}

\begin{column}{音の速さをこえた日}
音の速さに近づくと、飛行機の前に押し縮められた空気がたまり、機体が激しくゆれたり、かじがきかなくなったりしました。そのため、音の速さには「音の壁」があると言われました。1947年10月14日、アメリカのイェーガーが操縦するロケットの実験機ベル X-1 は、高さ約\qty{13}{km}で、マッハ約1.06（時速約\qty{1100}{km}）で水平に飛び、はじめて音の速さをこえました。マッハ1.06 のマッハ円すいは、半分の角が $\sin\alpha=1\div1.06$ より約\qty{71}{\degree}で、ほとんど平らな面に近い形です（⑥）。いまの旅客機は、マッハ0.8 くらい（熱-4 例題9）で、音の壁の手前を飛んでいます。
\end{column}

%% ======================== 章末問題 ========================
\bkendproblems
\bkproblemgroup{基本}{3}
\bkproblem \qty{512}{Hz}のおんさAと、高さのわからないおんさBを同時に鳴らすと、1秒に3回うなりました。Bに少しねん土をつけて（振動数を下げて）鳴らすと、うなりが1秒に5回にふえました。ねん土をつける前のBの振動数を求めましょう。
\bkpans{下げてうなりがふえたので、BはAより低い。$512-3=\qty{509}{Hz}$}
\bkproblem 2台のスピーカーから、同じ位相で\qty{1372}{Hz}の音が出ています（音の速さ\qty{343}{m/s}）。2台からの道のりの差が（1）\qty{0.375}{m}、（2）\qty{0.500}{m}の所では、強め合いますか、弱め合いますか。
\bkpans{$\lambda=\qty{0.250}{m}$。（1）1.5波長で弱め合う。（2）2波長で強め合う}
\bkproblem 地面の近くの\qty{0}{\celsius}の空気（音の速さ\qty{331}{m/s}）から、上の\qty{30}{\celsius}の空気（\qty{349}{m/s}）へ、音が入射角\qty{40}{\degree}で斜めに上がっていきます。屈折角を求めましょう。音は垂線に近づきますか、離れますか。
\bkpans{$\sin\theta_2=\sin40°\times349\div331\approx0.678$、約\qty{43}{\degree}。速い所へ入るので、垂線から離れる}
\bkproblem 汽笛（\qty{500}{Hz}）を鳴らす列車が\qty{30}{m/s}で走っています。線路のそばに止まって聞く人に、近づくときと遠ざかるときに聞こえる振動数を求めましょう（音の速さ\qty{343}{m/s}）。
\bkpans{近づくとき $343\div313\times500\approx\qty{548}{Hz}$。遠ざかるとき $343\div373\times500\approx\qty{460}{Hz}$}
\bkproblem \qty{1000}{Hz}の音が、幅\qty{0.50}{m}のすきまを通るとき、最初の静かな向きの角を求めましょう（音の速さ\qty{343}{m/s}）。波長 $0.60\times10^{-6}$\,m の光が、幅\qty{0.50}{mm}のすきまを通るときはどうですか。
\bkpans{音：$\sin\theta=0.343\div0.50\approx0.69$、約\qty{43}{\degree}。光：$\sin\theta=1.2\times10^{-3}$、約\qty{0.069}{\degree}}
\bkproblem 救急車（\qty{960}{Hz}）が\qty{15}{m/s}で、車に乗った人が\qty{20}{m/s}で、互いに近づいています。車の中の人に聞こえる振動数を求めましょう（音の速さ\qty{340}{m/s}）。
\bkpans{$v_{\mathrm s}=15$、$v_{\mathrm o}=-20$ で、$\tfrac{360}{325}\times960$、約\qty{1063}{Hz}}
\bkproblem \qty{600}{Hz}の音を出しながら、時速\qty{72}{km}でまっすぐ走るバイクがあります。バイクの前と後ろの、音の波長を求めましょう（音の速さ\qty{343}{m/s}）。
\bkpans{前 $(343-20)\div600\approx\qty{0.54}{m}$、後ろ $(343+20)\div600\approx\qty{0.61}{m}$}
\bkproblem 深さ\qty{2.0}{m}の所から深さ\qty{0.50}{m}の所へ、水の波が入射角\qty{20}{\degree}で入ります。屈折角を求めましょう（浅い水の波の速さ $\sqrt{gh}$）。
\bkpans{速さの比 $\sqrt{0.50\div2.0}=0.5$。$\sin\theta_2=\sin20°\times0.5\approx0.17$、約\qty{10}{\degree}}
\bkproblemgroup{チャレンジ}{4}
\bkproblem 10 GHz 帯のスピードガン（振動数 $10.525\times10^9$\,Hz）の前を、車が\qty{30}{m/s}で遠ざかっていきます。はね返った電波の振動数のずれを求めましょう（電波の速さ $3.00\times10^8$\,m/s）。
\bkpans{遠ざかるので $v=\qty{-30}{m/s}$ とする。$\Delta f\approx2vf/c$ より、約 $\qty{-2.1}{kHz}$（低くなる）}
\bkproblem 振動数\qty{4.0}{MHz}の超音波で、\qty{0.40}{m/s}で流れる血液をはかったら、ずれが\qty{1.0}{kHz}でした。超音波と血管のなす角を求めましょう（音の速さ\qty{1540}{m/s}）。
\bkpans{$\cos\theta=c\,\Delta f/(2fv)\approx0.48$、約\qty{61}{\degree}}
\bkproblem マッハ1.5の飛行機が、高さ\qty{9.0}{km}を水平に飛んでいます。音の速さを\qty{300}{m/s}として、マッハ円すいの半分の角を求め、真下の人にソニックブームが聞こえるのは、飛行機が真上を通ってから何秒後か求めましょう。
\bkpans{$\sin\alpha=1\div1.5$、$\alpha\approx41.8°$。先へ進む距離 $9000\div\tan41.8°\approx\qty{10100}{m}$。速さ\qty{450}{m/s}で、約\qty{22}{s}後}
\bkproblem 例題12のコウモリ（\qty{40}{kHz}、\qty{5.0}{m/s}）が、同じ向きに\qty{1.0}{m/s}で逃げるガを、まっすぐ追いかけています。ガが受けとる音の振動数と、コウモリが聞くこだまの振動数を、近似を使わずに求めましょう（音の速さ\qty{343}{m/s}）。近似の式 $2(v_{\text{コウモリ}}-v_{\text{ガ}})f/c$ と比べ、ガが止まっているとき（例題12）より、ずれが小さくなるわけを説明しましょう。
\bkpans{ガが受けとる：$\tfrac{343-1.0}{343-5.0}\times\num{40000}\approx\qty{40.47}{kHz}$。こだま：ガが音源として $\tfrac{343+5.0}{343+1.0}$ 倍にするので、約\qty{40.94}{kHz}。ずれは約\qty{0.94}{kHz}で、近似の\qty{0.93}{kHz}とほぼ同じ。ガも逃げているので、近づく速さが\qty{4.0}{m/s}と小さくなるから}

\begin{furikaeri}
\begin{checklist}
\item 近い高さの2つの音が重なるときの、うなりの回数を求められる。\bkfurref{ph:波-8:sec2}
\item ホイヘンスの原理で、反射・屈折の法則と回折を説明できる。\bkfurref{ph:波-8:sec1}
\item 2つの波源からの道のりの差で、強め合う所と弱め合う所を求められる。\bkfurref{ph:波-8:sec3}
\item 波面を数えてドップラー効果の式を導き、聞こえる高さを求められる。\bkfurref{ph:波-8:sec5}
\item 斜めのドップラー効果・反射板・衝撃波で、速さをはかるしくみを説明できる。\bkfurref{ph:波-8:sec7}
\end{checklist}
\end{furikaeri}
\bkfukidashi{delta}{救急車のサイレンが低くなるのは、波面の間かくがのびるからなんだね。スピードガンも血流の検査も、同じ「波面を数える」考えだったなんて！}
